% !TeX program = latexmk
\documentclass[10pt,a4paper]{amsart}
\usepackage{amssymb,mathtools}
\usepackage[margin=1in]{geometry}
\usepackage{extarrows}
\usepackage[all]{xy}
\usepackage{booktabs,tabularx,array,enumitem,microtype}
\usepackage[hidelinks]{hyperref}
\setlist{nosep,leftmargin=2.2em}
\allowdisplaybreaks[2]
\emergencystretch=1em

\theoremstyle{plain}
\newtheorem{theorem}{Theorem}[section]
\newtheorem{lemma}[theorem]{Lemma}
\newtheorem{proposition}[theorem]{Proposition}
\newtheorem{corollary}[theorem]{Corollary}
\theoremstyle{definition}
\newtheorem{definition}[theorem]{Definition}
\newtheorem{example}[theorem]{Example}
\newtheorem{conjecture}[theorem]{Conjecture}
\newtheorem{problem}[theorem]{Problem}
\theoremstyle{remark}
\newtheorem{remark}[theorem]{Remark}

\begin{document}

\title[Gabriel--Zisman Localizations and (Co)products]{Gabriel--Zisman Localizations, Products, Coproducts, and Product Categories}

\author{Chencheng Zhang}
\address{School of Mathematical Sciences, Shanghai Jiao Tong University, Shanghai 200240, P. R. China}
\email{zhangchencheng@sjtu.edu.cn}
\subjclass[2020]{Primary 18E35; Secondary 18A30, 18E05, 18N40}
\keywords{Gabriel--Zisman localization, product, coproduct, additive category, product category, calculus of fractions, model structure}
\thanks{The author was supported by the National Natural Science Foundation of China (No.~12131015).}
\begin{abstract}
  This paper studies (co)products and additive structures under Gabriel--Zisman localization $Q:\mathcal A\longrightarrow\mathcal A[S^{-1}]$.
  Examples show that localization may destroy existing (co)products or fail to preserve (co)products that remain, even when $Q$ is additive.
  For a set-indexed family $(\mathcal A_i,S_i)_{i\in I}$, necessary and sufficient conditions are given for the canonical functor $\big(\prod_{i\in I}\mathcal A_i\big)\big[\big(\prod_{i\in I}S_i\big)^{-1}\big]\longrightarrow\prod_{i\in I}\mathcal A_i[S_i^{-1}]$ to be full or faithful.
  These criteria yield compatibility with $I$-indexed (co)products for localizations admitting a calculus of fractions or arising from model categories, provided the relevant (co)products exist and the designated classes of morphisms are closed under them.
\end{abstract}

\maketitle

\section{Introduction}\label{sec:intro}
Given a category $\mathcal A$ and a class of morphisms $S\subseteq\mathsf{Mor}(\mathcal A)$, the Gabriel--Zisman localization $Q:\mathcal A\longrightarrow\mathcal A[S^{-1}]$ turns the morphisms in $S$ into isomorphisms, and every functor that sends $S$ to isomorphisms factors uniquely through $Q$ \cite[Chapter~I, \S~1, pp.~6--7]{GabrielZisman}.
This paper studies the relationships between Gabriel--Zisman localization and products, coproducts, and additive structures, with particular attention to two questions: if $\mathcal A$ has (co)products of a certain kind, does the localized category $\mathcal A[S^{-1}]$ still have them; if both $\mathcal{A}$ and $\mathcal A[S^{-1}]$ have (co)products of a certain kind, does the localization functor preserve them?

\begin{remark}
  All indexing sets and indexing categories in this paper are assumed to be small.
  If sizes need to be recorded, one may fix Grothendieck universes $\mathcal U\in\mathcal V$, use ``small'' to mean $\mathcal U$-small, and form the required categories and localizations in $\mathcal V$; universe decorations will be omitted below.
  For conventions concerning universes, see \cite[Expos\'e~I, \S~0 and Appendix, \S\S~1--5]{SGA4} and \cite[pp.~320--321, ``Conventions ensemblistes'']{CisinskiDerivable}; for size issues in localization, see \cite[Chapter~I, \S~1.3; Chapter~II, \S\S~2.1--2.2]{VerdierDerived} and \cite[pp.~425--426]{NandaUniverseFractions}.
\end{remark}

When $S$ contains all identity morphisms, $Q$ preserves finite (co)products whenever $\mathcal A$ has finite (co)products and $S$ is closed under finite (co)products.
Maltsiniotis proves the binary (co)product statement in \cite[Proposition~2.1.8, p.~75]{MaltsiniotisGrothendieck}; Tronin proves the finite-product statement in \cite[Theorem~1.3, p.~1526]{TroninProducts}, and the finite-coproduct statement follows by duality.
See Proposition~\ref{prop:finite-structure-localization} below for the precise form used here.
The proof descends the (co)product--diagonal adjunction to the localized categories, using the compatibility of localization with finite product categories.
To generalize this result to $I$-indexed (co)products,
one considers a family of nonempty categories $(\mathcal A_i)_{i\in I}$ and classes $S_i\subseteq\mathsf{Mor}(\mathcal A_i)$ containing all identity morphisms,
and therefore studies when the canonical functor is fully faithful:
\[
  \Theta:
  \bigg(\prod_{i\in I}\mathcal A_i\bigg)
  \bigg[\bigg(\prod_{i\in I}S_i\bigg)^{-1}\bigg]
  \longrightarrow
  \prod_{i\in I}\mathcal A_i[S_i^{-1}].
\]
It is known that $\Theta$ is an isomorphism of categories when $I$ is finite; see \cite[Lemma~2.1.7, pp.~74--75]{MaltsiniotisGrothendieck}, \cite[Theorem~1.2, pp.~1524--1525]{TroninProducts}, and Lemma~\ref{lem:localization-binary-product}.
For arbitrary $I$, $\Theta$ remains bijective on objects, but it need not be full or faithful.

\medskip

To state the main theorems, the notation $\ell$ and $d$ are introduced.
For every $i\in I$, let $\overline{S_i}$ denote the class of formal inverse arrows of the arrows in $S_i$, and let $\langle\mathsf{Mor}(\mathcal A_i)\sqcup\overline{S_i}\rangle$ denote the collection of all finite compositions of arrows in $\mathsf{Mor}(\mathcal A_i)\sqcup\overline{S_i}$.
For a composition of arrows $\mathfrak P$ in this collection, let $|\mathfrak P|$ be the number of arrows in $\mathfrak P$, let $[\mathfrak P]$ be the morphism that it represents in the localization, and put
\[
  \ell(\mathfrak P)
  =\min\bigl\{|\mathfrak P'|\bigm|[\mathfrak P']=[\mathfrak P]\bigr\}.
\]
For $\mathfrak P,\mathfrak Q \in \langle\mathsf{Mor}(\mathcal A_i)\sqcup\overline{S_i}\rangle$ with $[\mathfrak P]=[\mathfrak Q]$, let $\delta(\mathfrak P,\mathfrak Q)$ be the least number of elementary rewriting steps required to rewrite $\mathfrak P$ as $\mathfrak Q$; the elementary rewritings are given by the relations in the Gabriel--Zisman construction displayed in \eqref{eq:gz-arrow-composition-relations}.
For $L\in\mathbb N$, define
\[
  d_i(L)=\sup\left\{
  \delta(\mathfrak P,\mathfrak Q)\ \middle|\
  [\mathfrak P]=[\mathfrak Q],\
  |\mathfrak P|,|\mathfrak Q|\leq L
  \right\}\in\mathbb N\cup\{\infty\},
\]
where the supremum of the empty set is taken to be $0$.

\medskip

\noindent\textbf{Theorem.} \ Given a small set $I$, a family of nonempty categories $(\mathcal A_i)_{i\in I}$, and classes of morphisms $S_i\subseteq\mathsf{Mor}(\mathcal A_i)$ containing all identity morphisms, consider the canonical functor induced by localization
\[
  \Theta:
  \bigg(\prod_{i\in I}\mathcal A_i\bigg)
  \bigg[\bigg(\prod_{i\in I}S_i\bigg)^{-1}\bigg]
  \longrightarrow
  \prod_{i\in I}\mathcal A_i[S_i^{-1}].
\]
\begin{enumerate}[label=\textup{(\arabic*)}]
  \item (Theorem~\ref{thm:product-canonical-fullness}) $\Theta$ is full if and only if there exist a finite subset $F\subseteq I$ and $N\in\mathbb N$ such that $\ell(\mathfrak P)\leq N$ for every $i \notin F$ and every $\mathfrak P\in\langle\mathsf{Mor}(\mathcal A_i)\sqcup\overline{S_i}\rangle$.
  \item (Theorem~\ref{thm:product-canonical-faithfulness}) $\Theta$ is faithful if and only if, for every $L\in\mathbb N$, there exist a finite subset $F_L\subseteq I$ and $D_L\in\mathbb N$ such that $d_i(L)\leq D_L$ for every $i \notin F_L$.
\end{enumerate}

\medskip

Moreover, when $(\mathcal{A}_i, S_i)_{i \in I}$ are all identical, one obtains the following theorem.

\medskip

\noindent\textbf{Theorem.}  (Theorem~\ref{cor:fraction-calculus-preservation}) Let $\mathcal{A}$ be a category, assume that the class of morphisms $S \subseteq \mathsf{Mor}(\mathcal{A})$ contains all identity morphisms, and let $Q : \mathcal{A} \longrightarrow \mathcal{A}[S^{-1}]$ be the localization.
Suppose that the category $\mathcal{A}$ and the class of morphisms $S$ satisfy the following two conditions:
\begin{align*}
  N = \sup \big\{\ell(\mathfrak P) \mid \mathfrak P\in \langle\mathsf{Mor}(\mathcal A)\sqcup\overline{S}\rangle\big\} < \infty, \\
  \sup\big\{ \delta(\mathfrak P,\mathfrak Q) \mid [\mathfrak P]=[\mathfrak Q], \ \mathfrak P, \mathfrak Q\in \langle\mathsf{Mor}(\mathcal A)\sqcup\overline{S}\rangle, \ |\mathfrak P|, |\mathfrak Q| \leq N+1\big\} < \infty.
\end{align*}
Then the following statements hold:
\begin{enumerate}[label=\textup{(\arabic*)}]
  \item If $\mathcal{A}$ has $I$-indexed coproducts and $S$ is closed under $I$-indexed coproducts, then $\mathcal{A}[S^{-1}]$ also has $I$-indexed coproducts and $Q$ preserves them;
  \item If $\mathcal{A}$ has $I$-indexed products and $S$ is closed under $I$-indexed products, then $\mathcal{A}[S^{-1}]$ also has $I$-indexed products and $Q$ preserves them.
\end{enumerate}

\medskip

In particular, localizations arising from a calculus of fractions or a model structure satisfy the two conditions in Theorem~\ref{cor:fraction-calculus-preservation}.
See Corollary~\ref{thm:fraction-calculus-application} and Theorem~\ref{thm:model-category-application} for the corresponding results.
The coproduct statement under a calculus of left fractions was proved in \cite[Proposition~3.5.1]{KrauseLocalizationTheory}.
For ordinary model categories, the equivalence of the comparison functor for discrete diagram categories is also recorded in \cite[Section~2.5, equation~(2.7) and footnote~14, pp.~41--42]{RiehlHomotopicalCategories}.
The proofs below recover these conclusions from the two quantitative criteria above and use Egger's weakened model-category axioms, which require only finite products and finite coproducts.

\medskip

The paper is organized as follows.
Section~\ref{sec:preliminaries} recalls basic properties of the Gabriel--Zisman localization $Q : \mathcal{A} \longrightarrow \mathcal{A}[S^{-1}]$ and sufficient conditions under which a localization functor preserves finite (co)products and direct sums.
Section~\ref{sec:finite-counterexamples} gives examples in which $\mathcal{A}$ has finite (co)products but $\mathcal{A}[S^{-1}]$ does not, and examples in which both $\mathcal{A}$ and $\mathcal{A}[S^{-1}]$ have finite (co)products but $Q$ does not preserve them.
Under the assumptions that $\mathcal{A}$ and $\mathcal{A}[S^{-1}]$ are additive and that $Q$ is additive, Section~\ref{sec:additive-infinite-counterexamples} gives, respectively, an example in which $\mathcal{A}$ has arbitrary (co)products but $\mathcal{A}[S^{-1}]$ does not, and an example in which both $\mathcal{A}$ and $\mathcal{A}[S^{-1}]$ have arbitrary (co)products but $Q$ does not preserve them.
Section~\ref{sec:product-localization} establishes the necessary and sufficient conditions for $\Theta$ to be full and faithful, together with the corresponding preservation theorem.
Section~\ref{sec:discrete-applications} applies these results to calculus of fractions and model structures.
Section~\ref{sec:conjectures-and-problems} discusses the additivity of localization functors and the localization of general functor categories, and poses several questions.

\section{Preliminaries}\label{sec:preliminaries}

\subsection{Categorical notation}

Let $\mathcal{A}$ be a category.
An object $\bot$ is called an initial object of $\mathcal{A}$ if, for every $X \in \mathsf{Ob}(\mathcal{A})$, the class of morphisms $\mathrm{Hom}_\mathcal{A}(\bot, X)$ is a singleton; dually, an object $\top$ is called a terminal object of $\mathcal{A}$ if, for every $X \in \mathsf{Ob}(\mathcal{A})$, the class of morphisms $\mathrm{Hom}_\mathcal{A}(X, \top)$ is a singleton.
Fix an indexing set $I$.
A coproduct of a family of objects $(X_i)_{i \in I}$, if it exists, consists of an object $\coprod_{i \in I} X_i$ and a family of structure morphisms $e_i : X_i \longrightarrow \coprod_{i \in I} X_i$ such that, for every $Y \in \mathsf{Ob}(\mathcal{A})$, there is an isomorphism
\begin{equation}\label{eq_coprod}
  \mathrm{Hom}_\mathcal{A}\Big(\coprod_{i \in I} X_i, Y\Big) \xlongrightarrow{\sim} \prod_{i \in I} \mathrm{Hom}_\mathcal{A}(X_i, Y),\qquad f \longmapsto (f \circ e_i)_{i \in I}.
\end{equation}
Similarly, a product of $(X_i)_{i \in I}$, if it exists, consists of an object $\prod_{i \in I} X_i$ and a family of structure morphisms $p_i : \prod_{i \in I} X_i \longrightarrow  X_i$ such that, for every $Z \in \mathsf{Ob}(\mathcal{A})$, there is an isomorphism
\begin{equation}\label{eq_prod}
  \mathrm{Hom}_\mathcal{A}\Big(Z, \prod_{i \in I} X_i\Big) \xlongrightarrow{\sim} \prod_{i \in I} \mathrm{Hom}_\mathcal{A}(Z, X_i),\qquad g \longmapsto (p_i \circ g)_{i \in I}.
\end{equation}
The isomorphisms \eqref{eq_coprod} and \eqref{eq_prod} are also called the universal properties of the coproduct and product, respectively.
By convention, $\coprod_{i \in \varnothing} X_i$ is an initial object and $\prod_{i \in \varnothing} X_i$ is a terminal object.

\medskip

A category $\mathcal{A}$ is called preadditive if, for every $X,Y,Z \in \mathsf{Ob}(\mathcal{A})$, $\mathrm{Hom}_\mathcal{A}(X,Y)$ is an abelian group and composition is additive in both variables:
\[
  \mathrm{Hom}_\mathcal{A}(Y,Z) \times \mathrm{Hom}_\mathcal{A}(X,Y) \longrightarrow \mathrm{Hom}_\mathcal{A}(X,Z),\qquad (g,f) \longmapsto g \circ f.
\]
A category $\mathcal{A}$ is called additive if it is preadditive and has finite biproducts.
Equivalently, $\mathcal A$ is preadditive, finite coproducts and finite products exist, and, for every finite family of objects $(X_i)_{i \in I}$, the canonical morphism $\coprod_{i \in I}X_i\longrightarrow\prod_{i \in I}X_i$ is an isomorphism.
The finitely indexed coproduct and product are then denoted by $\bigoplus$.

Let $\mathcal{A}$ and $\mathcal{B}$ be additive categories.
A functor $F : \mathcal{A} \longrightarrow \mathcal{B}$ is called additive if, for every $X,Y \in \mathsf{Ob}(\mathcal{A})$, the following map is a homomorphism of additive groups:
\[
  \mathrm{Hom}_\mathcal{A}(X,Y) \longrightarrow \mathrm{Hom}_\mathcal{B}(FX,FY),\qquad f \longmapsto Ff.
\]
Equivalently, for every $X,Y \in \mathsf{Ob}(\mathcal{A})$, the universal properties of the biproduct in \eqref{eq_coprod} and \eqref{eq_prod} induce an isomorphism $F(X \oplus Y) \cong FX \oplus FY$.

\subsection{Gabriel--Zisman localization and saturated classes}

Given a category $\mathcal{A}$ and a class of morphisms $S$, the construction of the Gabriel--Zisman localization is recalled from \cite[Chapter~I, \S~1, pp.~6--7]{GabrielZisman}.
Forget the identity and composition laws on $\mathsf{Mor} (\mathcal{A})$.
A morphism is then regarded merely as an arrow; the source and target of this arrow remain those of the original morphism.
Let $\overline S = \{\overline s : Y \longrightarrow X\mid s : X \longrightarrow Y\text{ belongs to }S\}$ be a family of new arrows.
Form the disjoint union of arrows $\mathsf{Mor} (\mathcal{A}) \sqcup \overline S$.
Arrows $g$ and $f$ are called composable if the source of $g$ is the target of $f$.
Write their composition as $g \ast f$.

\begin{remark}
  In particular, the composite $g \circ f$ of arrows $f,g\in\mathsf{Mor}(\mathcal A)$ gives a composition of arrows $g \ast f$ in $\mathsf{Mor} (\mathcal{A}) \sqcup \overline S$; however, $g \circ f$ is itself an arrow in $\mathsf{Mor} (\mathcal{A}) \sqcup \overline S$, and it is not equal to the composition of arrows $g \ast f$.
\end{remark}
\noindent Let $\langle \mathsf{Mor} (\mathcal{A}) \sqcup \overline S \rangle$ denote the class of all finite compositions of arrows in $\mathsf{Mor} (\mathcal{A}) \sqcup \overline S$.
Consider the equivalence relation on $\langle \mathsf{Mor} (\mathcal{A}) \sqcup \overline S \rangle$ generated by the following rules:
\begin{enumerate}[label=\textup{(\arabic*)}]
  \item if $f,g \in \mathsf{Mor}(\mathcal{A})$ and $g$ and $f$ are composable, then $g \circ f$ and $g \ast f$ are equivalent;
  \item for every arrow $a : M \longrightarrow N$, the three compositions of arrows $\operatorname{Id}_N \ast a$, $a \ast \operatorname{Id}_M$, and $a$ are equivalent;
  \item if $s : X \longrightarrow Y$ is an arrow in $S$, then $s \ast \overline s$ is equivalent to $\operatorname{Id}_Y$, and $\overline s \ast s$ is equivalent to $\operatorname{Id}_X$;
  \item if the compositions of arrows $\mathfrak P$ and $\mathfrak Q$ are equivalent, then $a \ast \mathfrak P \ast b$ and $a \ast \mathfrak Q \ast b$ are equivalent whenever these compositions of arrows are defined.
\end{enumerate}

\medskip

Let $\langle \mathsf{Mor} (\mathcal{A}) \sqcup \overline S \rangle /_\sim$ be the quotient of $\langle \mathsf{Mor} (\mathcal{A}) \sqcup \overline S \rangle$ by this equivalence relation.
One obtains the following description of $\mathcal{A}[S^{-1}]$: its class of objects is $\mathsf{Ob}(\mathcal{A})$; the class of morphisms $\mathrm{Hom}_{\mathcal{A}[S^{-1}]}(X, Y)$ is the subclass of $\langle \mathsf{Mor} (\mathcal{A}) \sqcup \overline S \rangle/_\sim$ consisting precisely of compositions of arrows from $X$ to $Y$; the identity morphism of $X \in \mathsf{Ob}(\mathcal{A}[S^{-1}])$ is induced by the arrow $\operatorname{Id}_X$; composition of morphisms is induced by $\ast$.

\medskip

By construction, a morphism in $\mathcal{A}[S^{-1}]$ is written as $[\mathfrak P]$, where $\mathfrak P \in \langle \mathsf{Mor} (\mathcal{A}) \sqcup \overline S \rangle$ is a composition of arrows and $[-]$ denotes the equivalence class under the relation above.

\subsection{Basic facts about localization}

The following basic facts about localization will be used below.

\begin{lemma}\label{lem:nullary-preservation}
  Let $\mathcal{A}$ be a category and $S \subseteq \mathsf{Mor}(\mathcal{A})$ a class of morphisms.
  Let $Q:\mathcal A\longrightarrow \mathcal A[S^{-1}]$ be the localization.
  \begin{enumerate}[label=\textup{(\arabic*)}]
    \item if $\mathcal{A}$ has an initial object $\bot$ of $\mathcal A$, then $Q\bot$ is an initial object of $\mathcal A[S^{-1}]$;
    \item if $\mathcal{A}$ has a terminal object $\top$ of $\mathcal A$, then $Q\top$ is a terminal object of $\mathcal A[S^{-1}]$;
    \item if $\mathcal{A}$ has a zero object $0$ of $\mathcal A$, then $Q0$ is a zero object of $\mathcal A[S^{-1}]$.
  \end{enumerate}
\end{lemma}

\begin{proof}
  Suppose that $\bot$ is an initial object.
  For every $A\in\mathsf{Ob}(\mathcal A)$, let $u_A:\bot\longrightarrow A$ denote the unique morphism.
  First, one proves that $\mathrm{Hom}_{\mathcal A[S^{-1}]}(Q\bot,QA)$ is a singleton.
  Take any $[\mathfrak P] \in \mathrm{Hom}_{\mathcal A[S^{-1}]}(Q\bot,QA)$, where $\mathfrak P \in \langle \mathsf{Mor} (\mathcal{A}) \sqcup \overline S \rangle$.
  One may assume that $S$ contains all identity morphisms.
  By formally inserting identity arrows, $\mathfrak P$ is equivalent to a composition of arrows of the form
  \[
    \mathfrak P' = f_n \ast \overline {s_n} \ast f_{n-1} \ast \overline{s_{n-1}} \ast \cdots \ast f_1 \ast \overline{s_1} \ast f_0,
  \]
  where $f_0$ is a morphism from the initial object, written as $f_0 = u_X : \bot \longrightarrow X$.
  For $s_1 : Y \longrightarrow X$ and $f_1 : Y \longrightarrow Z$, one has $s_1 \circ u_Y = u_X$ and $f_1 \circ u_Y = u_Z$.
  Hence $[f_1 \ast \overline{s_1} \ast f_0] = [u_Z]$.
  Inductively, $\mathfrak P'$ is equivalent to $u_A$.
  Thus $\mathrm{Hom}_{\mathcal A[S^{-1}]}(Q\bot,QA)$ is a singleton, so $Q \bot$ is initial.

  \medskip

  The argument for a terminal object is similar.
  Consequently, localization also preserves zero objects.
\end{proof}

\begin{lemma}\label{lem_loc_uni}
  Let $\mathcal{A}$ be a category and $S \subseteq \mathsf{Mor}(\mathcal{A})$ a class of morphisms.
  Then localization is unique up to an isomorphism of categories, rather than merely an equivalence of categories.
  Equivalently, if $Q_1:\mathcal A\longrightarrow \mathcal{L}_1$ and $Q_2:\mathcal A\longrightarrow \mathcal{L}_2$ are both localizations defined by the universal property, then there is an isomorphism $\mathcal{L}_1 \cong \mathcal{L}_2$.
\end{lemma}

\begin{proof}
  Let $\widetilde {Q_2} : \mathcal{L}_1 \longrightarrow \mathcal{L}_2$ be the functor induced by the universal property of the localization $Q_1$, and let $\overline{Q_1} : \mathcal{L}_2 \longrightarrow \mathcal{L}_1$ be the functor induced by the universal property of $Q_2$.
  By construction, $Q_1 = \overline{Q_1} \circ \widetilde{Q_2} \circ Q_1$ and $Q_2 = \widetilde{Q_2} \circ \overline{Q_1} \circ Q_2$.
  The universal property of localization gives $\operatorname{Id}_{\mathcal{L}_1} = \overline{Q_1} \circ \widetilde{Q_2}$ and $\operatorname{Id}_{\mathcal{L}_2} = \widetilde{Q_2} \circ \overline{Q_1}$.
  This constructs an isomorphism between the categories $\mathcal{L}_1$ and $\mathcal{L}_2$.
\end{proof}

\begin{lemma}\label{lem:localization-binary-product}
  Let $I$ be a finite indexing set.
  For every $i \in I$, let $\mathcal{A}_i$ be a category and $S_i \subseteq \mathsf{Mor}(\mathcal{A}_i)$ a class of morphisms.
  Assume that $S_i$ contains all identity morphisms.
  Let $Q_i : \mathcal{A}_i \longrightarrow \mathcal{A}_i[S_i^{-1}]$ be the localization.
  Then the following functor induced by the universal property of localization is an isomorphism of categories:
  \[
    \widetilde {\prod_{i \in I} Q_i} :
    \bigg(\prod_{i \in I} \mathcal{A}_i \bigg)\bigg[\big(\prod_{i \in I}S_i\big)^{-1}\bigg]
    \longrightarrow\prod_{i \in I}\mathcal{A}_i[S_i^{-1}].
  \]
\end{lemma}

\begin{proof}
  See \cite[Lemma~2.1.7, pp.~74--75]{MaltsiniotisGrothendieck}.
\end{proof}

The binary product and coproduct assertions in the following proposition are proved in \cite[Proposition~2.1.8, p.~75]{MaltsiniotisGrothendieck}.
The product assertion also follows from \cite[Theorem~1.3, p.~1526]{TroninProducts}, while \cite[Theorem~3.1, pp.~1535--1536]{TroninProducts} gives the corresponding preadditive conclusion.
Since the descent of the coproduct--diagonal adjunction in the finite case provides the basic idea for the main theorems below, the details are included here.

\begin{proposition}\label{prop:finite-structure-localization}
  Let $\mathcal{A}$ be a category and $S \subseteq \mathsf{Mor}(\mathcal{A})$ a class of morphisms.
  Let $Q:\mathcal A\longrightarrow \mathcal A[S^{-1}]$ be the localization.
  Assume further that $S$ contains all identity morphisms: this does not change the localization, but makes the statement more convenient.
  \begin{enumerate}[label=\textup{(\arabic*)}]
    \item If $\mathcal{A}$ has finite coproducts and $S$ is closed under finite coproducts, then $\mathcal{A}[S^{-1}]$ has finite coproducts and $Q$ preserves them;
    \item If $\mathcal{A}$ has finite products and $S$ is closed under finite products, then $\mathcal{A}[S^{-1}]$ has finite products and $Q$ preserves them;
    \item If $\mathcal{A}$ is an additive category and $S$ is closed under finite biproducts, then $\mathcal{A}[S^{-1}]$ is an additive category and $Q$ preserves finite biproducts.
  \end{enumerate}
\end{proposition}

\begin{proof}
  It suffices first to prove (1); assertion (2) is dual.
  (3) follows from (1) and (2).

  \medskip

  By Lemma~\ref{lem:nullary-preservation}, it remains only to prove the following: if $\mathcal{A}$ has finite coproducts and $S$ is closed under finite coproducts, then $\mathcal{A}[S^{-1}]$ has binary coproducts and $Q$ preserves them.
  The binary coproduct functor $\coprod : \mathcal{A} \times \mathcal{A} \longrightarrow \mathcal{A}$ and the diagonal functor $\Delta : \mathcal{A} \longrightarrow \mathcal{A} \times \mathcal{A}$ are respectively left and right adjoint.
  In particular, the right adjoint $\Delta$ sends $S$ to $S \times S$; since $S$ is closed under binary coproducts, the left adjoint $\coprod$ sends $S \times S$ to $S$.
  By \cite[Lemma~1.1.6, p.~9]{KrauseHomologicalTheory}, or by the primary source \cite[Chapter~I, \S~1, Lemma~1.2, p.~7]{GabrielZisman}, the adjunction $\coprod \dashv \Delta$ therefore induces an adjunction on the localized categories
  \[
    \xymatrix@C=30mm@R=10mm{
    \mathcal A\times\mathcal A
    \ar@<0.8ex>[r]^{\coprod}
    \ar[d]_{Q\times Q}
    &
    \mathcal A
    \ar@<0.8ex>[l]^{\Delta}
    \ar[d]^{Q}
    \\
    \mathcal A[S^{-1}]\times\mathcal A[S^{-1}]
    \ar@<0.8ex>[r]^-{\coprod_S}
    &
    \mathcal A[S^{-1}]
    \ar@<0.8ex>[l]^-{\Delta_S}
    } \qquad \xymatrix@R=10mm{
    Q \circ \coprod = \coprod _S \circ (Q \times Q),\\
    (Q \times Q) \circ \Delta = \Delta_S \circ Q.
    }
  \]
  By Lemma~\ref{lem:localization-binary-product}, one may regard $Q \times Q$ as the localization functor of $\mathcal{A} \times \mathcal{A}$ at $S \times S$.
  It is clear that $\Delta_S$ is still the diagonal functor and has a left adjoint.
  Consequently, $\mathcal{A}[S^{-1}]$ has binary coproducts and hence finite coproducts.
  The commutative diagram shows that the localization $Q$ preserves coproducts.
\end{proof}

\section{Localization and (co)products}\label{sec:finite-counterexamples}

\subsection{Examples in which \texorpdfstring{$\mathcal{A}$}{} has finite (co)products but \texorpdfstring{$\mathcal{A}[S^{-1}]$}{} does not}

Consider a localization $Q:\mathcal A\longrightarrow\mathcal A[S^{-1}]$.
The following example shows that even if $\mathcal{A}$ has finite (co)products, $\mathcal{A}[S^{-1}]$ need not have finite (co)products.

\begin{example}\label{ex_seven_element_lattice}
  Let $\mathcal A$ be the poset category determined by the following Hasse diagram:
  \[
    \vcenter{\hbox{$\xymatrix@C=8mm@R=4mm{
      &A\ar[dr]^{\gamma}&&B\ar[dr]^{\eta}&\\
      \bot\ar[ur]^{\alpha}\ar[dr]_{\beta}&&
      J\ar[ur]^{\varepsilon}\ar[dr]_{\zeta}&&\top\\
      &X\ar[ur]_{\delta}&&Y\ar[ur]_{\theta}&
      }$}}.
  \]
  This is a finite lattice, so $\mathcal A$ has all finite products and finite coproducts.
  Let $S=\{s\} = \{\varepsilon\circ\gamma\}$.

  \medskip

  It will be shown that $QX \coprod QX$ and $QY \prod QY$ do not exist in $\mathcal A[S^{-1}]$.
\end{example}

\begin{proof}
  Take any morphism $[\mathfrak P]$ in $\mathcal{A}[S^{-1}]$, where $\mathfrak P \in \langle\mathsf{Mor}(\mathcal A)\sqcup\overline S\rangle$ is a composition of arrows.
  If two arrows from $\overline {S}$ occur in $\mathfrak P$, then between them there is an arrow or a composition of arrows from $A$ to $B$; hence $\mathfrak P$ has the form
  \[
    \cdots \ast \overline s \ast s \ast \overline s \ast \cdots \quad \text{or} \quad \cdots \ast \overline s \ast \varepsilon \ast \gamma \ast \overline s \ast \cdots.
  \]
  Since $\overline s \ast s \ast \overline s$ and $\overline s \ast \varepsilon \ast \gamma \ast \overline s$ are both equivalent to $\overline s$, one may assume that $\overline s$ occurs at most once in $\mathfrak P$.
  Therefore, for every $M, N \in \mathcal{A}$,
  \begin{equation}\label{eq_hom_bound}
    |\mathrm{Hom}_{\mathcal{A}[S^{-1}]}(QM,QN)|
    \leq \underset{\substack{\text{number of morphisms}\\\text{containing no }\overline s}}{\underbracket{|\mathrm{Hom}_{\mathcal{A}}(M,N)|}}
    \quad + \underset{\substack{\text{number of morphisms containing}\\\text{exactly one }\overline s}}{\underbracket{|\mathrm{Hom}_{\mathcal{A}}(M,B)| \cdot |\mathrm{Hom}_{\mathcal{A}}(A,N)|}}
    \leq 2.
  \end{equation}
  Next, one proves that $\mathrm{Hom}_{\mathcal A[S^{-1}]}(QX,QY)$ contains two distinct morphisms.
  Together with \eqref{eq_hom_bound}, this will show that its cardinality is $2$.
  Let $k$ be a field.
  Consider the functor $F : \mathcal{A} \to \mathrm{Vect}_k$ given by
  \[
    F :\quad
    \vcenter{\hbox{$\xymatrix@C=6mm@R=4mm{
      &A\ar[dr]^{\gamma}&&B\ar[dr]^{\eta}&\\
      \bot\ar[ur]^{\alpha}\ar[dr]_{\beta}&&
      J\ar[ur]^{\varepsilon}\ar[dr]_{\zeta}&&\top\\
      &X\ar[ur]_{\delta}&&Y\ar[ur]_{\theta}&
      }$}} \quad
    \longrightarrow \quad
    \vcenter{\hbox{$\xymatrix@C=6mm@R=4mm{
      &k\ar[dr]^{e_1}&&k\ar[dr]^{0}&\\
      0\ar[ur]^{0}\ar[dr]_{0}&&
      k^2\ar[ur]^{p_1}\ar[dr]_{p_2}&&0\\
      &k\ar[ur]_{e_2}&&k\ar[ur]_{0}&
      }$}}.
  \]
  By construction, $Fs = \operatorname{Id}_k$.
  The universal property of localization gives a unique functor
  \[
    \widetilde F:\mathcal A[S^{-1}]
    \longrightarrow\operatorname{Vect}_k,
    \qquad \widetilde F\circ Q=F.
  \]
  The morphisms $[\zeta\ast\delta]$ and $[\zeta\ast\gamma\ast\overline s\ast \varepsilon\ast\delta]$ are distinct in $\mathrm{Hom}_{\mathcal{A}[S^{-1}]}(QX, QY)$.
  By construction,
  \begin{align*}
    \widetilde F([\zeta\ast\delta])
     & =F(\zeta) \circ F(\delta)
    =p_2\circ e_2
    =\operatorname{Id}_k,                                                             \\
    \widetilde F([\zeta\ast\gamma\ast\overline s\ast \varepsilon\ast\delta])
     & =F(\zeta) \circ F(\gamma) \circ F(s)^{-1} \circ F(\varepsilon) \circ F(\delta)
    =p_2\circ e_1\circ \operatorname{Id}_k^{-1}\circ p_1\circ e_2
    =0.
  \end{align*}
  Therefore $\widetilde F([\zeta\ast\delta]) \neq \widetilde F([\zeta\ast\gamma\ast\overline s\ast \varepsilon\ast\delta])$, and hence $[\zeta\ast\delta] \neq [\zeta\ast\gamma\ast\overline s\ast \varepsilon\ast\delta]$.
  This proves that $\left| \mathrm{Hom}_{\mathcal A[S^{-1}]}(QX,QY) \right|=2$.

  \medskip

  It remains to show that $QX \coprod QX$ does not exist.
  Otherwise, \eqref{eq_coprod} would give
  \[
    \big|\mathrm{Hom}_{\mathcal A[S^{-1}]}(QX\coprod QX,QY)\big| = \big|\mathrm{Hom}_{\mathcal A[S^{-1}]}(QX,QY)\big|^2 = 4.
  \]
  This contradicts \eqref{eq_hom_bound}.
  If $QY\prod QY$ existed, then \eqref{eq_prod} would similarly give
  \[
    \big|\mathrm{Hom}_{\mathcal A[S^{-1}]}(QX,QY\prod QY)\big|
    =\big|\mathrm{Hom}_{\mathcal A[S^{-1}]}(QX,QY)\big|^2=4,
  \]
  again contradicting \eqref{eq_hom_bound}.
\end{proof}

The following example shows that a Gabriel--Zisman localization of an additive category need not be additive.
A complete proof valid over an arbitrary field is given below; in fact, the resulting localization is not even preadditive.

\begin{example}\label{ex:vector-space-nonpreadditive-localization}
  Let $\mathcal A$ be the category of finite-dimensional $k$-vector spaces, let $S = \{0 \xlongrightarrow s k\}$, and let $Q : \mathcal{A} \longrightarrow \mathcal{A}[S^{-1}]$ be the localization.
  Then $\mathcal A[S^{-1}]$ is not an additive category.
\end{example}

\begin{proof}
  By Lemma~\ref{lem:nullary-preservation}, $Q0$ is a zero object of $\mathcal A[S^{-1}]$.
  Since $Qs:Q0\longrightarrow Qk$ is an isomorphism, $Qk$ is also a zero object.
  Moreover $Q$ is full.
  Indeed, $Q(s)^{-1}$ is represented by $k \longrightarrow 0$.
  It follows that for each $X,Y \in \mathcal{A}$, $\mathrm{Hom}_{\mathcal{A}[S^{-1}]}(QX,QY)$ is a quotient of $\mathrm{Hom}_{\mathcal{A}}(X,Y)$.

  \medskip

  Consider the functor
  \[
    \Phi:\mathcal A\longrightarrow\mathcal A, \qquad
    \Phi(V)=\bigwedge\nolimits^2V \oplus\bigl(\bigwedge\nolimits^2V\bigr)\otimes_kV,\quad
    \Phi(f)=\bigwedge\nolimits^2f \oplus\bigl(\bigwedge\nolimits^2f\bigr)\otimes_kf,
  \]
  Since both $\Phi k$ and $\Phi 0$ are zero objects, $\Phi s$ is an isomorphism.
  The universal property of localization gives a functor $\widetilde\Phi:\mathcal{A}[S^{-1}]\longrightarrow\mathcal A$ satisfying $\widetilde\Phi\circ Q=\Phi$.
  Fix an isomorphism $\bigwedge^2k^2\cong k$.
  For every $g\in\operatorname{GL}_2(k)$, a direct calculation gives $\Phi(g)=\det(g)\oplus\bigl(\det(g) \cdot g\bigr)$.
  Therefore, if $g,h\in\operatorname{GL}_2(k)$ and $[g]=[h]$, comparing the first components in this formula gives $\det(g)=\det(h)$, and then comparing the second components gives $g=h$.
  It follows that
  \[
    \operatorname{GL}_2(k)\longrightarrow
    \operatorname{End}_{\mathcal A [S^{-1}]}(Qk^2),
    \qquad g\longmapsto[g],
  \]
  is injective.
  Moreover, $\Phi(\operatorname{Id}_{k^2})\neq\Phi(0)$, so $[\operatorname{Id}_{k^2}]\neq[0]$.

  \medskip

  Now suppose, for a contradiction, that $\mathcal A [S^{-1}]$ is preadditive.
  Since $Qk$ is a zero object, every rank $1$ linear map $f:k^2\longrightarrow k^2$ satisfies $[f] = [0]$.
  On the other hand, every rank $2$ linear map is an isomorphism.
  Since $Q$ is full, every element of $\operatorname{End}_{\mathcal A [S^{-1}]}(Qk^2)$ has the form $[f]$.
  Hence the ring $D = \operatorname{End}_{\mathcal A [S^{-1}]}(Qk^2)$ is nonzero and every nonzero element is invertible; thus $D$ is a division ring, so $1_D$ has at most two square roots.
  If $\operatorname{char}(k)\neq2$, then $\pm \begin{psmallmatrix}-1&0\\0&1\end{psmallmatrix}$ give a contradiction; if $\operatorname{char}(k)=2$, then $\begin{psmallmatrix}1&1\\0&1\end{psmallmatrix}$ and $\begin{psmallmatrix}1&0\\1&1\end{psmallmatrix}$ give a contradiction.

  \medskip

  Thus $\mathcal A[S^{-1}]$ is not preadditive and hence is not additive.
\end{proof}

\subsection{Examples in which \texorpdfstring{$\mathcal{A}$}{} and \texorpdfstring{$\mathcal{A}[S^{-1}]$}{} have finite (co)products that are not preserved by \texorpdfstring{$Q$}{}}

Consider a localization $Q:\mathcal A\longrightarrow\mathcal A[S^{-1}]$.
The following example shows that even if both $\mathcal{A}$ and $\mathcal{A}[S^{-1}]$ have finite (co)products, $Q$ need not preserve them.

\begin{example}\label{ex:diamond-lattice}
  Let $\mathcal A$ be the poset category determined by the following diamond lattice:
  \[
    \xymatrix@C=8mm@R=4mm{
    &A\ar[dr]^{\gamma}&\\
    \bot\ar[ur]^{\alpha}\ar[dr]_{\beta}&&\top\\
    &B\ar[ur]_{\delta}&
    }
  \]
  Let $S=\{\alpha\}$ and let $Q:\mathcal A\longrightarrow\mathcal A[S^{-1}]$ be the localization.
  Both $\mathcal A$ and $\mathcal A[S^{-1}]$ have finite coproducts, but $Q$ does not preserve binary coproducts.


\end{example}

\begin{proof}
  Let $\mathbf 3=(0<1<2)$ be the poset category determined by the three-element chain.
  Define a functor $F:\mathcal A\longrightarrow\mathbf 3$ by
  \[
    F(\bot)=F(A)=0,
    \qquad F(B)=1,
    \qquad F(\top)=2.
  \]
  Then $F(\alpha)=\operatorname{Id}_0$.
  By the universal property of localization, $F$ factors through $Q$ as $\widetilde F:\mathcal A[S^{-1}]\longrightarrow\mathbf 3$, with $\widetilde F\circ Q=F$.

  \medskip

  The functor $\widetilde F$ is an equivalence of categories.
  In the other direction, define a functor $G:\mathbf 3\longrightarrow \mathcal A[S^{-1}]$.
  On objects and morphisms it is given by
  \[
    G : \quad 0 \longrightarrow 1 \longrightarrow 2 \quad \longmapsto \quad
    QA \xlongrightarrow{[\beta \ast \overline \alpha]} QB \xlongrightarrow{[\delta]} Q\top.
  \]
  Here $[\delta\ast\beta\ast\overline\alpha] = [\gamma\ast\alpha\ast\overline\alpha]= [\gamma]$.
  This shows that $G$ is indeed a functor.
  A direct calculation gives $\widetilde FG=\operatorname{Id}_{\mathbf 3}$.
  Define a natural isomorphism $\eta:\operatorname{Id}_{\mathcal A[S^{-1}]}\cong G\widetilde F$; its components are the vertical arrows in the following diagram:
  \[
    \xymatrix@C=13mm@R=8mm{
    \operatorname{Id}&
    Q\bot\ar@<.55ex>[r]^{[\alpha]}\ar[d]_{[\alpha]}&
    QA\ar@<.55ex>[l]^{[\overline\alpha]}
    \ar[r]^{[\beta\ast\overline\alpha]}
    \ar[d]^{[\operatorname{Id}_A]}&
    QB\ar[r]^{[\delta]}\ar[d]^{[\operatorname{Id}_B]}&
    Q\top\ar[d]^{[\operatorname{Id}_{\top}]}\\
    G\widetilde F&
    QA\ar@<.55ex>[r]^{[\operatorname{Id}_A]}&
    QA\ar@<.55ex>[l]^{[\operatorname{Id}_A]}
    \ar[r]_{[\beta\ast\overline\alpha]}&
    QB\ar[r]_{[\delta]}&
    Q\top.
    }
  \]
  Since $\mathbf 3$ is a bounded lattice, it has finite products and finite coproducts, and therefore so does $\mathcal{A}[S^{-1}]$.
  But $Q$ does not preserve binary coproducts.
  Indeed,
  \[
    Q (A \coprod B) = Q \top \ncong QB \cong QA \coprod QB.
  \]
  This gives an example in which both $\mathcal{A}$ and $\mathcal{A}[S^{-1}]$ have finite coproducts, but $Q$ does not preserve binary coproducts.
\end{proof}

Passing to opposite categories gives an example in which both $\mathcal{A}$ and $\mathcal{A}[S^{-1}]$ have finite products, but $Q$ does not preserve binary products.

\section{Additive localizations and infinite (co)products}\label{sec:additive-infinite-counterexamples}

\subsection{\texorpdfstring{An example in which $Q$ is additive, $\mathcal{A}$ has arbitrary coproducts, but $\mathcal{A}[S^{-1}]$ does not}{An example in which Q is additive, A has arbitrary coproducts, but its localization does not}}

Let $Q : \mathcal{A} \longrightarrow \mathcal{A}[S^{-1}]$ be a localization.
Assume that $\mathcal{A}$ and $\mathcal{A}[S^{-1}]$ are additive categories and that $Q$ is an additive functor.
Even if $\mathcal{A}$ has arbitrary coproducts, $\mathcal{A}[S^{-1}]$ need not have arbitrary coproducts.

\begin{example}\label{ex:additive-localization-without-countable-coproduct}
  Let $\mathcal A$ be the category of free abelian groups; its objects have the form $\coprod_{i\in I}\mathbb Z$, where $I$ is a set, and its morphisms are group homomorphisms.
  Let
  \[
    S=\left\{s:F\longrightarrow G\ \middle|\
    s\text{ is injective, and there exists $n\geq0$ such that }2^n\operatorname{coker}(s)=0\right\}.
  \]
  Then $\mathcal A$ has coproducts, $\mathcal A[S^{-1}]$ is an additive category, and the localization $Q:\mathcal A\longrightarrow\mathcal A[S^{-1}]$ is additive; however, the coproduct of the countable family $(Q\mathbb Z)_{n\geq1}$ does not exist in $\mathcal A[S^{-1}]$.
\end{example}

\begin{proof}
  It is clear that $S$ is closed under finite biproducts.
  By Proposition~\ref{prop:finite-structure-localization}(3), $\mathcal A[S^{-1}]$ is an additive category and $Q$ is an additive functor.

  \medskip

  Let $R=\mathbb Z[1/2]$ and consider the tensor functor
  \[
    T=-\otimes_{\mathbb Z}R:
    \mathcal A\longrightarrow R\text{-}\mathsf{Mod}.
  \]
  Since $R$ is a flat $\mathbb Z$-module, the tensor functor is exact.
  For every $S$-morphism $s:F\longrightarrow G$, there is a short exact sequence
  \[
    0\longrightarrow F\otimes_{\mathbb Z}R
    \xlongrightarrow{T(s)}G\otimes_{\mathbb Z}R
    \longrightarrow\operatorname{coker}(s)\otimes_{\mathbb Z}R
    \longrightarrow0.
  \]
  Since $2$ is invertible in $R$ and $\operatorname{coker}(s)$ is annihilated by a power of $2$, one has $\operatorname{coker}(s)\otimes_{\mathbb Z}R=0$, so $T(s)$ is an isomorphism.
  By the universal property of localization, there is a unique functor $\widetilde T:\mathcal A[S^{-1}]\longrightarrow R\text{-}\mathsf{Mod}$ such that $T=\widetilde T\circ Q$.

  \medskip

  It remains to prove that the coproduct of $(Q\mathbb Z)_{n\geq1}$ does not exist.
  Suppose, for a contradiction, that an object $QC$ together with structure morphisms $(e_n:Q\mathbb Z\longrightarrow QC)_{n \geq 1}$ forms this coproduct.
  For each $n\geq1$, let $x_n=\widetilde T(e_n)(1)\in C\otimes_{\mathbb Z}R$.
  Regard $C$ as a subgroup of $C\otimes_{\mathbb Z}R = \bigcup_{d\geq0}2^{-d}C$.
  One may therefore choose $d_n\geq0$ such that $2^{d_n}x_n\in C$.
  For every $m\geq0$, define the multiplication map
  \[
    \mu_m:\mathbb Z\longrightarrow\mathbb Z,
    \qquad a\longmapsto2^ma.
  \]
  By construction, $\mu_m\in S$.
  Thus $[\overline{\mu_m}] : Q\mathbb Z \longrightarrow Q\mathbb Z$ is a well-defined morphism in $\mathcal{A}[S^{-1}]$.
  Applying the universal property of the coproduct to the family of morphisms $([\overline{\mu_{d_n+n}}])_{n\geq1}$, there is a unique morphism $u:QC\longrightarrow Q\mathbb Z$ such that $u\circ e_n=[\overline{\mu_{d_n+n}}]$ for all $n \geq 1$.

  Choose a finite arrow representation $u=[\mathfrak P]$.
  Suppose that all inverse arrows occurring in $\mathfrak P$ are $(\overline{s_i}:G_i\longrightarrow F_i)_{1 \leq i \leq r}$, where $s_i:F_i\longrightarrow G_i$ belongs to $S$.
  For every $i$, choose $N_i\geq0$ such that $2^{N_i}\operatorname{coker}(s_i)=0$.
  Then $T(s_i)^{-1}(G_i)\subseteq2^{-N_i}F_i$.
  Indeed, for every $y\in G_i$, one can choose $z\in F_i$ such that $s_i(z)=2^{N_i}y$, and hence $T(s_i)^{-1}(y)=2^{-N_i}z$.
  An ordinary arrow does not increase the denominator, whereas an inverse arrow $\overline{s_i}$ multiplies the denominator by at most $2^{N_i}$.
  Composing successively along the arrows in $\mathfrak P$ therefore gives
  \[
    \widetilde T(u)(C)\subseteq2^{-N}\mathbb Z , \qquad (N =N_1+\cdots+N_r).
  \]
  If no inverse arrow occurs in $\mathfrak P$, take $N=0$ in this formula.
  Since $2^{d_n}x_n\in C$, one has
  \[
    2^{d_n}\widetilde T(u)(x_n)
    =\widetilde T(u)(2^{d_n}x_n)
    \in2^{-N}\mathbb Z.
  \]
  On the other hand, $u\circ e_n=[\overline{\mu_{d_n+n}}]$ gives
  \[
    \widetilde T(u)(x_n)
    =\widetilde T(u\circ e_n)(1)
    =\widetilde T([\overline{\mu_{d_n+n}}])(1)
    =2^{-(d_n+n)}.
  \]
  Choose $n>N$.
  Then $2^{-(d_n+n)}\notin2^{-(d_n+N)}\mathbb Z$, a contradiction.
\end{proof}

Passing to opposite categories in the preceding example shows that $\mathcal{A}$ and $\mathcal{A}[S^{-1}]$ may both be additive categories and $Q$ additive, while $\mathcal{A}$ has arbitrary products but $\mathcal{A}[S^{-1}]$ does not.

\subsection{\texorpdfstring{An example in which $Q$ is additive and $\mathcal{A}$ and $\mathcal{A}[S^{-1}]$ have arbitrary products that are not preserved by $Q$}{An example in which Q is additive and A and its localization have arbitrary products that are not preserved by Q}}

\begin{example}\label{ex:torsion-complexes}
  Let $\mathsf{TorsAb}$ denote the category of torsion abelian groups.
  An object $G \in \mathsf{Ab}$ is called a torsion abelian group if every $g \in G$ has finite order.
  Let $\mathcal A=\operatorname{Ch}(\mathsf{TorsAb})$ be the category of cochain complexes in $\mathsf{TorsAb}$.
  Let $S$ be the class of all quasi-isomorphisms in $\mathcal{A}$, and let $Q:\mathcal A\longrightarrow\mathcal A[S^{-1}]$ be the corresponding additive localization.
  Then both $\mathcal A$ and $\mathcal A[S^{-1}]$ have arbitrary products, but $Q$ does not preserve countable products.
\end{example}

\begin{proof}
  Let $t(G)$ denote the torsion subgroup of an abelian group $G$, namely the subgroup consisting of all elements of finite order.
  For any indexing set $I$ and any family of torsion abelian groups $(A_i)_{i\in I}$, their product in $\mathsf{TorsAb}$ is $t\left(\prod_{i\in I}A_i\right)$, where $\prod_{i\in I}A_i$ is the product in $\mathsf{Ab}$.
  Indeed, for any torsion abelian group $t(G)$ and any family of morphisms $(f_i:t(G)\longrightarrow A_i)_{i\in I}$, there is a unique morphism $f:t(G)\longrightarrow\prod_{i\in I}A_i$ such that $f_i=p_i\circ f$ for all $i\in I$, where $p_i:\prod_{i\in I}A_i\longrightarrow A_i$ is the structure morphism.
  In particular, the image of $f$ lies in $t\left(\prod_{i\in I}A_i\right)$, which shows that $t\left(\prod_{i\in I}A_i\right)$ is the product in $\mathsf{TorsAb}$.

  \medskip

  Fix a prime $p$.
  For every $n\geq1$, consider the epimorphism
  \[
    f_n:\mathbb Z/p^n\mathbb Z\longrightarrow\mathbb Z/p\mathbb Z,
    \qquad  a + p^n \mathbb Z \longmapsto a + p \mathbb Z .
  \]
  The family of morphisms $(f_n)_{n\geq1}$ induces the following epimorphism in $\mathsf{Ab}$:
  \[
    \varphi_0 : \prod\nolimits_{n\geq1}^{\mathsf{Ab}}\mathbb Z/p^n\mathbb Z
    \longrightarrow
    \prod\nolimits_{n\geq1}^{\mathsf{Ab}}\mathbb Z/p\mathbb Z,
    \qquad
    (x_n)_{n\geq1}\longmapsto(f_n(x_n))_{n\geq1}.
  \]
  Let $\varphi$ be the restriction of $\varphi_0$ to the torsion subgroup, namely
  \[
    \varphi:
    t\left(\prod_{n\geq1}\mathbb Z/p^n\mathbb Z\right)
    \longrightarrow
    \prod_{n\geq1}\mathbb Z/p\mathbb Z,
    \qquad
    (x_n)_{n\geq1}\longmapsto(f_n(x_n))_{n\geq1}.
  \]
  The morphism $\varphi$ is not surjective.
  Indeed, the element $(\overline1,\overline1,\ldots)$ is not in the image of $\varphi$.
  For any $(x_n)_{n \geq 1} \in \varphi_0 ^{-1}((\overline1,\overline1,\ldots))$, the equality $f_n(x_n)=\overline1$ implies that $x_n$ has order $p^n$.
  Hence no nonzero integer annihilates all components of $(x_n)_{n\geq1}$.
  Thus $(x_n)_{n \geq 1} \notin t\left(\prod_{n\geq1}\mathbb Z/p^n\mathbb Z\right)$, and $\varphi$ is not surjective.

  \medskip

  For every $n\geq1$, regard the short exact sequence
  \[
    0 \longrightarrow \ker(f_n) \longrightarrow \mathbb Z/p^n\mathbb Z \xlongrightarrow{f_n} \mathbb Z/p\mathbb Z \longrightarrow 0
  \]
  as an acyclic complex $K_n$ concentrated in degrees $0,1,2$.
  Let $K=\prod_{n\geq1}K_n$ be the product in $\mathcal A$.
  Since products are computed degreewise, the differential of $K$ from degree $1$ to degree $2$ is precisely $\varphi$.
  Because $\varphi$ is not surjective, $H^2(K)=\operatorname{coker}_{\mathsf{TorsAb}}(\varphi)\neq0$.
  Thus $K$ is not acyclic.

  The derived category is the localization $Q : \mathcal{A} \longrightarrow \mathcal A[S^{-1}]=D(\mathsf{TorsAb})$.
  Here $Q$ is the composite of two additive functors $\mathcal A\longrightarrow K(\mathsf{TorsAb}) \longrightarrow D(\mathsf{TorsAb})$.
  The category $\mathsf{TorsAb}$ is a Grothendieck abelian category\footnote{The category $\mathsf{TorsAb}$ is a Serre subcategory of $\mathsf{Ab}$ and is closed under arbitrary coproducts, so the assertion follows from \cite[Proposition~2.2.16 and its proof, pp.~36--37]{KrauseHomologicalTheory}. Taking the generator $\mathbb Z$ of $\mathsf{Ab}$ in the construction given there yields the generator $\coprod_{m\geq1}\mathbb Z/m\mathbb Z$ of $\mathsf{TorsAb}$.}.
  By Serp\'e's theorem on $K$-injective resolutions \cite[Theorem~3.13]{SerpeKInjective}, every unbounded complex has a $K$-injective resolution.
  The degreewise product of a family of $K$-injective complexes $(J_\lambda)_{\lambda\in\Lambda}$ is again $K$-injective: for every acyclic complex $X$, there is a canonical isomorphism
  \[
    \operatorname{Hom}^{\bullet}\!\left(X,\prod_{\lambda\in\Lambda}J_\lambda\right)
    \cong
    \prod_{\lambda\in\Lambda}\operatorname{Hom}^{\bullet}(X,J_\lambda),
  \]
  and arbitrary products of abelian groups preserve exact sequences.
  Therefore $D(\mathsf{TorsAb})$ has arbitrary products, which may be computed as degreewise products of $K$-injective resolutions.

  \medskip

  Finally, every $K_n$ is acyclic, so $\prod_{n\geq1}QK_n\cong0$.
  But $H^2(K)\neq0$, so $QK\not\cong0$.
  Hence the canonical morphism induced by the product projections
  \[
    Q\left(\prod_{n\geq1}K_n\right)
    =QK
    \longrightarrow
    \prod_{n\geq1}QK_n
  \]
  is not an isomorphism.
  This shows that $Q$ does not preserve countable products.
\end{proof}

Passing to opposite categories in the preceding example gives the following example: $\mathcal{A}$ and $\mathcal{A}[S^{-1}]$ are additive categories with arbitrary coproducts, while $Q$ is additive but does not preserve countable coproducts.

\section{Localization of product categories}\label{sec:product-localization}

The study of whether localization preserves finite (co)products used the isomorphism in Lemma~\ref{lem:localization-binary-product}.
Likewise, when studying whether localization preserves arbitrary (co)products, the main obstruction is the canonical functor
\[
  \Theta : \widetilde {\prod_{i \in I} Q_i} :
  \bigg(\prod_{i \in I} \mathcal{A}_i \bigg)\bigg[\big(\prod_{i \in I}S_i\big)^{-1}\bigg]
  \longrightarrow\prod_{i \in I}\mathcal{A}_i[S_i^{-1}],
\]
where $I$ is an arbitrary indexing set, not necessarily finite.
Henceforth every $\mathcal A_i$ is assumed to be nonempty, and $S_i$ is assumed to contain all identity morphisms of $\mathcal A_i$.
If some $\mathcal A_i$ is empty, then the source and target of $\Theta$ are both empty categories; if $I=\varnothing$, then both are terminal categories.
In either case, $\Theta$ is an isomorphism of categories.

\subsection{A necessary and sufficient condition for \texorpdfstring{$\Theta$}{Theta} to be full}
For every $i\in I$, define
\[
  \ell:
  \langle\mathsf{Mor}(\mathcal A_i)\sqcup\overline{S_i}\rangle
  \longrightarrow\mathbb N,
  \qquad
  \ell(\mathfrak P)
  =
  \min\bigl\{|\mathfrak P'|\bigm|[\mathfrak P']=[\mathfrak P]\bigr\},
\]
where $|\mathfrak P|$ denotes the number of arrows in $\mathfrak P$.

\begin{theorem}\label{thm:product-canonical-fullness}
  The canonical functor
  \[
    \Theta : \widetilde {\prod_{i \in I} Q_i} :
    \bigg(\prod_{i \in I} \mathcal{A}_i \bigg)\bigg[\big(\prod_{i \in I}S_i\big)^{-1}\bigg]
    \longrightarrow\prod_{i \in I}\mathcal{A}_i[S_i^{-1}]
  \]
  is full if and only if there exist a finite subset $F\subseteq I$ and $N\in\mathbb N$ such that $\ell(\mathfrak P)\leq N$ for every $i\notin F$ and every $\mathfrak P\in \langle\mathsf{Mor}(\mathcal A_i)\sqcup\overline{S_i}\rangle$.
\end{theorem}

\begin{proof}
  Suppose first that $\Theta$ is full.
  Assume for contradiction that the stated condition fails, i.e.,
  \begin{equation}\label{eq:fullness-bound-negation}
    \forall E\subseteq I\text{ finite},\quad
    \forall n\in\mathbb N,\quad
    \exists i\in I\setminus E,\quad
    \exists\mathfrak P_i\in
    \langle\mathsf{Mor}(\mathcal A_i)\sqcup\overline{S_i}\rangle,
    \quad
    \ell(\mathfrak P_i)>n.
  \end{equation}
  Under this assumption, $I$ must be infinite; fix a sequence $\{i_0, i_1, \ldots\}$.
  Set $E_1 = \{i_0\}$ and choose the corresponding $i_1 \in I \setminus E_1$ by \eqref{eq:fullness-bound-negation}.
  For $n \geq 1$, define $E_n=\{i_0,\ldots,i_{n-1}\}$ inductively.
  One obtains pairwise distinct elements $i_n\in I$ and compositions of arrows
  \begin{equation}\label{eq:fullness-unbounded-compositions}
    \mathfrak P_{i_n}:X_{i_n}\longrightarrow Y_{i_n},
    \qquad
    \ell(\mathfrak P_{i_n})>n
    \qquad(n\geq 1).
  \end{equation}
  For $i\notin\{i_n\mid n\in\mathbb N\}$, let $\mathfrak P_i$ be any identity arrow.
  This gives $([\mathfrak P_i])_{i\in I}\in \prod_{i\in I} \operatorname{Hom}_{\mathcal A_i[S_i^{-1}]} (Q_iX_i,Q_iY_i)$.
  Since $\Theta$ is full, there is a composition of arrows $\mathfrak R$ in $\langle\mathsf{Mor}(\prod_{i\in I}\mathcal A_i)\sqcup\overline{\prod_{i\in I}S_i}\rangle$ for which the following diagram holds:
  \[
    \xymatrix@C=20mm@R=2mm{
    \operatorname{Hom}_{
    (\prod_i\mathcal A_i)[(\prod_iS_i)^{-1}]}
    ((X_i)_i,(Y_i)_i)
    \ar[r]^-{\Theta}
    &
    \prod_i\operatorname{Hom}_{\mathcal A_i[S_i^{-1}]}
    (Q_iX_i,Q_iY_i)
    \\
    [\mathfrak R]\ar@{|->}[r]
    &
    ([\mathfrak R_i])_{i\in I}
    =([\mathfrak P_i])_{i\in I}.
    }
  \]
  Here $\mathfrak R_i$ is the $i$th coordinate of $\mathfrak R$.
  Let $m=|\mathfrak R|$.
  Then $\ell(\mathfrak P_i) \leq |\mathfrak R_i| \leq |\mathfrak R| = m$ holds.
  Taking $n>m$ in \eqref{eq:fullness-unbounded-compositions} gives $n<\ell(\mathfrak P_{i_n})\leq m<n$, a contradiction.

  \medskip

  Conversely, suppose that there exist a finite subset $F\subseteq I$ and $N\in\mathbb N$ such that
  \begin{equation}\label{eq:fullness-assumed-bound}
    \forall i\notin F,\quad
    \forall \mathfrak P\in
    \langle\mathsf{Mor}(\mathcal A_i)\sqcup\overline{S_i}\rangle, \quad \text{one has}
    \quad
    \ell(\mathfrak P)\leq N.
  \end{equation}
  It remains to prove that $\Theta$ is full.
  Take any morphism $([\mathfrak P_i])_{i\in I}:(Q_iX_i)_{i\in I}\longrightarrow(Q_iY_i)_{i\in I}$ in $\prod_{i \in I}\mathcal{A}_i[S_i^{-1}]$.
  For $i\notin F$, use \eqref{eq:fullness-assumed-bound} to choose $\mathfrak R_i$ such that $[\mathfrak R_i]=[\mathfrak P_i]$ and $|\mathfrak R_i|=\ell(\mathfrak P_i)\leq N$.
  For $i\in F$, set $\mathfrak R_i=\mathfrak P_i$.
  Let
  \begin{equation}\label{eq:fullness-global-composition-bound}
    M=\max\bigl(\{N\}\cup
    \{|\mathfrak R_i|\mid i\in F\}\bigr).
  \end{equation}
  Then $|\mathfrak R_i|\leq M$ for every $i\in I$.
  By inserting identity arrows, write each $\mathfrak R_i$ as an equivalent composition of arrows
  \begin{equation}\label{eq:common-alternating-representatives}
    \mathfrak R_i'
    =f_{i,M}\ast\overline{s_{i,M}}\ast f_{i,M-1}\ast
    \cdots\ast f_{i,1}\ast\overline{s_{i,1}}\ast f_{i,0},\qquad
    [\mathfrak R_i']=[\mathfrak R_i]=[\mathfrak P_i],\qquad
    |\mathfrak R_i'| = 2M+1.
  \end{equation}
  where all $f_{i,j}\in\mathsf{Mor}(\mathcal A_i)$ and $s_{i,j}\in S_i$.
  For $0\leq j\leq M$ and $1\leq j\leq M$, respectively, put
  \[
    f_j=(f_{i,j})_{i\in I}\in
    \mathsf{Mor}\left(\prod_{i\in I}\mathcal A_i\right),
    \qquad
    s_j=(s_{i,j})_{i\in I}\in\prod_{i\in I}S_i.
  \]
  The compositions of arrows in \eqref{eq:common-alternating-representatives} assemble coordinatewise into the composition of arrows $\mathfrak R=f_M\ast\overline{s_M}\ast f_{M-1}\ast\cdots\ast f_1\ast\overline{s_1}\ast f_0$.
  Thus
  \[
    \Theta ([\mathfrak R]) = ([\mathfrak R_i'])_{i\in I} =([\mathfrak P_i])_{i\in I}.
  \]
  Therefore $\Theta$ is surjective on every Hom set; that is, $\Theta$ is full.
\end{proof}

The following example shows that $\Theta$ need not be full.

\begin{example}\label{ex:product-localization-not-full}
  Let $I=\mathbb N_{>0}$, and for each $i\in I$ take the same one-object category $\mathcal C$, whose object is $\ast$ and whose morphisms form the monoid $(\mathbb N,+)$.
  Let $S=\{0,1\}\subseteq\mathbb N$, where $0 = \operatorname{Id}_\ast$ is the identity morphism.
  The following canonical functor is not full:
  \[
    \Theta:(\mathcal C^{\mathbb N_{>0}})[(S^{\mathbb N_{>0}})^{-1}]
    \longrightarrow(\mathcal C[S^{-1}])^{\mathbb N_{>0}}.
  \]
\end{example}

\begin{proof}
  One has $\operatorname{End}_{\mathcal C[S^{-1}]}(\ast)\cong(\mathbb Z,+)$.
  Indeed, take any $[\mathfrak P] \in \mathcal{C}[S^{-1}]$.
  Up to equivalence, one may assume that the only inverse arrow in $\mathfrak P$ is that of $1$, denoted by $-1$.
  It is clear that $n \ast 1$ is equivalent to $1 \ast n$, and hence $\overline 1 \ast n$ is equivalent to $n \ast \overline 1$.
  Also, $\overline 1 \ast 1$ is equivalent to $0$.
  Therefore $\mathfrak P$ is equivalent either to an arrow in $\mathsf{Mor}(\mathcal C)$ or to a composition of arrows of the form $\underbracket{\overline 1\ast\cdots\ast\overline 1}_{q\text{ copies}}$.
  It is straightforward to verify that
  \[
    \mathbb N \sqcup \big\{\underbracket{\overline 1\ast\cdots\ast\overline 1}_{q\text{ copies}}\mid q \geq 1\big\} \longrightarrow \mathbb Z, \qquad n \longmapsto n, \quad \underbracket{\overline 1\ast\cdots\ast\overline 1}_{q\text{ copies}} \longmapsto (-q)
  \]
  is a bijection of sets, under which $\ast$ corresponds to addition.
  Thus $\mathrm{End}_{\mathcal{C}[S^{-1}]}(\ast) \cong \mathbb Z$.

  \medskip

  For $n\geq1$, let $[\mathfrak P_n]=-n\in\mathbb Z$.
  There is no $\mathfrak R$ such that $\Theta ([\mathfrak R])=([\mathfrak P_n])_{n \geq 1}$.
  Otherwise, let $m$ be the number of inverse arrows in $\mathfrak R$.
  Then every component of $\mathfrak R$ contains at most $m$ inverse arrows; but $[\mathfrak R]$ has value $-(m+1)$ in its $(m+1)$st component, a contradiction.
  Hence $\Theta$ is not full.
\end{proof}

\subsection{A necessary and sufficient condition for \texorpdfstring{$\Theta$}{Theta} to be faithful}

For every $i\in I$, fix the following five types of elementary rewriting.
Here $f:X\longrightarrow Y$ and $g:Y\longrightarrow Z$ are arrows in $\mathsf{Mor}(\mathcal A_i)$, $a:M\longrightarrow N$ is an arrow in $\mathsf{Mor}(\mathcal A_i)\sqcup\overline{S_i}$, $s:X\longrightarrow Y$ is an arrow in $S_i$, and $\mathfrak U,\mathfrak V$ are compositions of arrows for which the following expressions are defined.
Either $\mathfrak U$ or $\mathfrak V$ may be omitted.
\begin{equation}\label{eq:gz-arrow-composition-relations}
  \begin{aligned}
    \mathfrak U\ast(g\ast f)\ast\mathfrak V
     & \longleftrightarrow
    \mathfrak U\ast(g\circ f)\ast\mathfrak V,          \\
    \mathfrak U\ast(\operatorname{Id}_N\ast a)\ast\mathfrak V
     & \longleftrightarrow
    \mathfrak U\ast a\ast\mathfrak V,                  \\
    \mathfrak U\ast(a\ast\operatorname{Id}_M)\ast\mathfrak V
     & \longleftrightarrow
    \mathfrak U\ast a\ast\mathfrak V,                  \\
    \mathfrak U\ast(s\ast\overline s)\ast\mathfrak V
     & \longleftrightarrow
    \mathfrak U\ast\operatorname{Id}_Y\ast\mathfrak V, \\
    \mathfrak U\ast(\overline s\ast s)\ast\mathfrak V
     & \longleftrightarrow
    \mathfrak U\ast\operatorname{Id}_X\ast\mathfrak V.
  \end{aligned}
\end{equation}
For any two equivalent compositions of arrows $\mathfrak P$ and $\mathfrak Q$, define
\[
  \delta(\mathfrak P,\mathfrak Q)
  =
  \min\left\{
  m\in\mathbb N\ \middle|\
  \begin{array}{c}
    \mathfrak P=\mathfrak R_0
    \longleftrightarrow\mathfrak R_1
    \longleftrightarrow\cdots
    \longleftrightarrow\mathfrak R_m=\mathfrak Q, \\
    \text{each step is one elementary rewriting from \eqref{eq:gz-arrow-composition-relations}}
  \end{array}
  \right\}.
\]
For $L\in\mathbb N$, define
\begin{equation}\label{eq:def-di-L}
  d_i(L) = \sup\left\{ \delta(\mathfrak P,\mathfrak Q) \mid [\mathfrak P]=[\mathfrak Q],\quad |\mathfrak P|\leq L,\quad|\mathfrak Q|\leq L \right\} \in\mathbb N\cup\{\infty\}.
\end{equation}
Here the supremum of the empty set is taken to be $0$.
For $L\geq1$, the set in \eqref{eq:def-di-L} is nonempty because $\delta(\operatorname{Id}_X,\operatorname{Id}_X)=0$.

If $\mathfrak{P}$ and $\mathfrak Q$ are related by at most $n$ elementary rewriting steps, write $\mathfrak{P}\xlongleftrightarrow{\leq n} \mathfrak{Q}$.

\begin{theorem}\label{thm:product-canonical-faithfulness}
  The canonical functor
  \[
    \Theta:\widetilde{\prod_{i\in I}Q_i}:
    \bigg(\prod_{i\in I}\mathcal A_i\bigg)
    \bigg[\bigg(\prod_{i\in I}S_i\bigg)^{-1}\bigg]
    \longrightarrow
    \prod_{i\in I}\mathcal A_i[S_i^{-1}]
  \]
  is faithful if and only if, for every $L\in\mathbb N$, there exist a finite subset $F_L\subseteq I$ and $D_L\in\mathbb N$ such that $d_i(L)\leq D_L$ for every $i\notin F_L$.
\end{theorem}

\begin{proof}
  Suppose first that $\Theta$ is faithful.
  If the stated boundedness condition fails, then there is an $L\in\mathbb N$ such that
  \begin{equation}\label{eq:faithfulness-bound-negation}
    \forall E\subseteq I\text{ finite},\ \
    \forall n\in\mathbb N,\ \
    \exists i\in I\setminus E,\ \
    \exists\mathfrak P_i,\mathfrak Q_i,\ \
    |\mathfrak P_i|,|\mathfrak Q_i|\leq L,\ \
    [\mathfrak P_i]=[\mathfrak Q_i],\ \
    \delta(\mathfrak P_i,\mathfrak Q_i)>n.
  \end{equation}
  One may assume that $L\geq1$.
  Under this assumption, $I$ must be infinite; fix a sequence $\{i_0, i_1, \ldots\}$.
  Set $E_1 = \{i_0\}$ and choose the corresponding $i_1 \in I \setminus E_1$ by \eqref{eq:faithfulness-bound-negation}.
  For $n \geq 1$, define $E_n=\{i_0,\ldots,i_{n-1}\}$ inductively.
  One obtains pairwise distinct elements $i_n\in I$ and corresponding pairs of compositions of arrows $(\mathfrak P_{i_n},\mathfrak Q_{i_n})$.

  For a composition of arrows $\mathfrak R$, let $\operatorname{sgn}(\mathfrak R)$ be the sign string obtained by recording its arrows in their composition order: an arrow in $\mathsf{Mor}(\mathcal A_i)$ is denoted by $+$ and an arrow in $\overline{S_i}$ by $-$.
  Since there are only finitely many sign strings of length at most $L$, one may pass to a subsequence of $\{i_k\}_{k \geq 0}$ and reindex it so that fixed sign strings $\sigma$ and $\tau$ satisfy
  \begin{equation}\label{eq:fixed-sign-pairs}
    \operatorname{sgn}(\mathfrak P_{i_n})=\sigma,\quad
    \operatorname{sgn}(\mathfrak Q_{i_n})=\tau,\quad
    [\mathfrak P_{i_n}]=[\mathfrak Q_{i_n}],\quad
    \delta(\mathfrak P_{i_n},\mathfrak Q_{i_n})>n
    \qquad(n\geq1),
  \end{equation}
  For $i\notin\{i_n\mid n\geq1\}$, choose an object $X_i\in\mathcal A_i$, and let $\mathfrak P_i$ and $\mathfrak Q_i$ be the compositions of arrows built from $\operatorname{Id}_{X_i}$ and $\overline{\operatorname{Id}_{X_i}}$ whose sign strings are $\sigma$ and $\tau$, respectively.
  Let $\mathfrak P = (\mathfrak P_i)_{i\in I}$ and $\mathfrak Q = (\mathfrak Q_i)_{i\in I}$.
  By construction,
  \begin{equation}\label{eq:same-theta-image}
    \Theta[\mathfrak P]
    =([\mathfrak P_i])_{i\in I}
    =([\mathfrak Q_i])_{i\in I}
    =\Theta[\mathfrak Q].
  \end{equation}
  By assumption, $\Theta$ is faithful.
  Therefore $[\mathfrak P]=[\mathfrak Q]$.
  Hence there is a chain of elementary rewritings $\mathfrak P=\mathfrak R_0 \longleftrightarrow\cdots \longleftrightarrow\mathfrak R_m=\mathfrak Q$.
  Projecting this chain to the $i_n$th coordinate and using \eqref{eq:fixed-sign-pairs}, one obtains $n<\delta(\mathfrak P_{i_n},\mathfrak Q_{i_n})\leq m$ for every $n \geq 1$.
  Taking $n=m+1$ gives a contradiction.
  This completes the proof in this direction.

  \medskip

  Conversely, assume that the stated boundedness condition holds.
  Take morphisms $[\mathfrak P],[\mathfrak Q]$ with the same source and target such that $\Theta[\mathfrak P]=\Theta[\mathfrak Q]$.
  Let $\mathfrak P_i,\mathfrak Q_i$ be the $i$th components of $\mathfrak P, \mathfrak Q$, respectively.
  Put $L=\max\{|\mathfrak P|,|\mathfrak Q|\}$.
  Then $|\mathfrak P_i|,|\mathfrak Q_i|\leq L$ and $[\mathfrak P_i]=[\mathfrak Q_i]$ for every $i\in I$.
  Choose $F_L,D_L$ as in the hypothesis, and let
  \begin{equation}\label{eq:uniform-rewriting-bound}
    D=\max\bigl(\{D_L\}\cup
    \{\delta(\mathfrak P_i,\mathfrak Q_i)\mid i\in F_L\}\bigr).
  \end{equation}
  Thus, for each $i\in I$, one may choose a chain of $m_i\leq D$ elementary rewritings
  \begin{equation}\label{eq:coordinate-rewriting-chain}
    \mathfrak P_i=\mathfrak R_{i,0}
    \longleftrightarrow\cdots
    \longleftrightarrow\mathfrak R_{i,m_i}=\mathfrak Q_i.
  \end{equation}
  An elementary rewriting increases the number of arrows by at most $1$, so $|\mathfrak R_{i,j}|\leq L+D$.
  For the $j$th step, let $\nu_{i,j},\varepsilon_{i,j},p_{i,j}$ denote, respectively, the row number of the relation in \eqref{eq:gz-arrow-composition-relations}, the direction in which it is used, and the position where it is applied.
  Record
  \[
    \kappa_i=\left(
    m_i;\
    (\nu_{i,j},\varepsilon_{i,j},p_{i,j})_{1\leq j\leq m_i};\
    (\operatorname{sgn}(\mathfrak R_{i,j}))_{0\leq j\leq m_i}
    \right).
  \]
  Since $m_i\leq D$, there are only five types of elementary relation, and the positions and sign strings all have length at most $L+D$, the invariant $\kappa_i$ takes only finitely many values.
  Partition $I$ into finitely many blocks according to these values, $I=I_1\amalg\cdots\amalg I_r$, where $i,i'\in I_t$ if and only if $\kappa_i=\kappa_{i'}$.
  Put $\mathcal B_t=\prod_{i\in I_t}\mathcal A_i$ and $T_t=\prod_{i\in I_t}S_i$.
  Within each $I_t$, every coordinate uses a relation from the same row, in the same direction, and with the same sign string at every step.
  Therefore the coordinatewise chains in \eqref{eq:coordinate-rewriting-chain} assemble into a chain of elementary rewritings in $\mathcal B_t$:
  \[
    \mathfrak P^{(t)}=\mathfrak R^{(t)}_0
    \longleftrightarrow\cdots
    \longleftrightarrow\mathfrak R^{(t)}_{m_t}=\mathfrak Q^{(t)},
    \qquad m_t=m_i\quad(i\in I_t),
  \]
  where $\mathfrak P^{(t)},\mathfrak Q^{(t)}$ are the projections of $\mathfrak P,\mathfrak Q$ to $I_t$, respectively.
  Hence $[\mathfrak P^{(t)}]=[\mathfrak Q^{(t)}]$ as morphisms in $\mathcal B_t[T_t^{-1}]$.

  From $\prod_{i\in I}\mathcal A_i\cong\prod_{t=1}^r\mathcal B_t$ and Lemma~\ref{lem:localization-binary-product}, one obtains an isomorphism of categories
  \[
    \Phi:
    \left(\prod_{i\in I}\mathcal A_i\right)
    \left[\left(\prod_{i\in I}S_i\right)^{-1}\right]
    \xlongrightarrow{\sim}
    \prod_{t=1}^r\mathcal B_t[T_t^{-1}].
  \]
  Thus $\Phi[\mathfrak P] =\bigl([\mathfrak P^{(t)}]\bigr)_{t=1}^r =\bigl([\mathfrak Q^{(t)}]\bigr)_{t=1}^r =\Phi[\mathfrak Q]$, hence $[\mathfrak P]=[\mathfrak Q]$.
  Therefore $\Theta$ is faithful.
\end{proof}

\begin{proposition}\label{prop:faithfulness-under-fullness}
  Suppose that $\Theta$ is full.
  By Theorem~\ref{thm:product-canonical-fullness}, there exist a finite subset $F_0\subseteq I$ and an integer $N_0\geq1$ such that
  \[
    \ell(\mathfrak P)\leq N_0
    \qquad
    \bigl(i\notin F_0,\
    \mathfrak P\in
    \langle\mathsf{Mor}(\mathcal A_i)\sqcup\overline{S_i}\rangle\bigr).
  \]
  Then $\Theta$ is faithful if and only if there exist a finite subset $F\subseteq I$ and $D\in\mathbb N$ such that $d_i(N_0+1)\leq D$ for every $i\notin F$.
\end{proposition}

\begin{proof}
  If $\Theta$ is faithful, take $L=N_0+1$ in Theorem~\ref{thm:product-canonical-faithfulness} and enlarge $F_L$ by $F_0$ to obtain the stated condition.

  \medskip

  Conversely, assume that the stated condition holds.
  Replacing $F$ by $F\cup F_0$, one may assume that $F_0\subseteq F$.
  Fix $i\notin F$ and take a composition of arrows $\mathfrak P=a_m\ast\cdots\ast a_1$ in $\langle\mathsf{Mor}(\mathcal A_i)\sqcup\overline{S_i}\rangle$.
  Set $\mathfrak R_1=a_1$, and use the length bound $N_0$ to choose successively representatives satisfying $[\mathfrak R_k]=[a_k\ast\mathfrak R_{k-1}]$ and $|\mathfrak R_k|\leq N_0$, for $2\leq k\leq m$.
  Since $|a_k\ast\mathfrak R_{k-1}|\leq N_0+1$, one has
  \[
    \delta(a_k\ast\mathfrak R_{k-1},\mathfrak R_k)
    \leq d_i(N_0+1)\leq D.
  \]
  If $m\geq2$, then for $k=2,\ldots,m$, replace $a_k\ast\mathfrak R_{k-1}$ by $\mathfrak R_k$ using at most $D$ elementary rewriting steps.
  This gives
  \[
    \mathfrak P
    =a_m\ast\cdots\ast a_2\ast\mathfrak R_1\xlongleftrightarrow{\leq D}
    a_m\ast\cdots\ast a_3\ast\mathfrak R_2\xlongleftrightarrow{\leq D}\cdots
    \xlongleftrightarrow{\leq D}
    \mathfrak R_m.
  \]
  There are $m-1$ replacements, each using at most $D$ elementary rewriting steps.
  If $m=1$, then $\mathfrak P=\mathfrak R_1$ and no rewriting is needed.
  Thus in all cases
  \[
    [\mathfrak R_m]=[\mathfrak P],
    \qquad |\mathfrak R_m|\leq N_0,
    \qquad
    \delta(\mathfrak P,\mathfrak R_m)\leq(m-1)D.
  \]
  If another composition of arrows $\mathfrak Q=b_n\ast\cdots\ast b_1$ in $\langle\mathsf{Mor}(\mathcal A_i)\sqcup\overline{S_i}\rangle$ satisfies $[\mathfrak P]=[\mathfrak Q]$, applying the same construction to $\mathfrak Q$ gives a composition of arrows $\mathfrak T_n$ satisfying
  \[
    [\mathfrak T_n]=[\mathfrak Q],
    \qquad |\mathfrak T_n|\leq N_0,
    \qquad
    \delta(\mathfrak Q,\mathfrak T_n)\leq(n-1)D.
  \]
  Since $[\mathfrak R_m]=[\mathfrak T_n]$ and $|\mathfrak R_m|,|\mathfrak T_n|\leq N_0$, one has $\delta(\mathfrak R_m,\mathfrak T_n)\leq d_i(N_0+1)\leq D$.
  Hence
  \[
    \delta(\mathfrak P,\mathfrak Q)
    \leq
    \delta(\mathfrak P,\mathfrak R_m)
    +\delta(\mathfrak R_m,\mathfrak T_n)
    +\delta(\mathfrak T_n,\mathfrak Q)\leq(m-1)D+D+(n-1)D.
  \]
  Thus $d_i(L)\leq(2L-1)D$ for every $L\geq1$ and every $i\notin F$.
  By Theorem~\ref{thm:product-canonical-faithfulness}, $\Theta$ is faithful.
\end{proof}

\begin{example}\label{ex:product-localization-not-faithful}
  Let $\mathcal C$ be the full subcategory of $\mathbb F_2$-vector spaces whose objects are $\{\mathbb F_2^d\}_{d \geq 0}$.
  For $d\geq0$, let
  \[
    j_d:\mathbb F_2^d\longrightarrow\mathbb F_2^{d+1},
    \qquad
    (x_1,\ldots,x_d) \longmapsto (x_1,\ldots,x_d,0),
  \]
  and take $S=\{j_d\mid d\geq0\} \cup \{\text{all isomorphisms}\}$.
  The following canonical functor is not faithful:
  \[
    \Theta:
    (\mathcal C^{\mathbb N_{>0}})
    [(S^{\mathbb N_{>0}})^{-1}]
    \longrightarrow
    (\mathcal C[S^{-1}])^{\mathbb N_{>0}}.
  \]
\end{example}

\begin{proof}
  Let $X=(\mathbb F_2^n)_{n\geq1}$, and take the morphisms $u=(\operatorname{Id}_{\mathbb F_2^n})_{n\geq1}$ and $v=(0_{n\times n})_{n\geq1}$.
  It will be shown that $\Theta[u]=\Theta[v]$ but $[u]\ne[v]$.

  \medskip

  On the one hand, for every $n\geq1$, the zero morphism $z_n:0\longrightarrow\mathbb F_2^n$ factors as
  \[
    0=\mathbb F_2^0
    \xrightarrow{j_0}\mathbb F_2
    \xrightarrow{j_1}\mathbb F_2^2
    \longrightarrow\cdots
    \xrightarrow{j_{n-1}}\mathbb F_2^n.
  \]
  Hence $[z_n]$ is invertible in $\mathcal C[S^{-1}]$.
  Cancelling $[z_n]$ from $\operatorname{Id}_{\mathbb F_2^n}\circ z_n =z_n =0_{n\times n}\circ z_n$ gives $[\operatorname{Id}_{\mathbb F_2^n}] =[0_{n\times n}]$.
  Therefore
  \[
    \Theta[u]
    =([\operatorname{Id}_{\mathbb F_2^n}])_{n\geq1}
    =([0_{n\times n}])_{n\geq1}
    =\Theta[v].
  \]

  \medskip

  On the other hand, define a relation on each Hom set of $\mathcal C^{\mathbb N_{>0}}$ by
  \[
    (f_n)_{n\geq1}\sim(g_n)_{n\geq1}
    \quad\Longleftrightarrow\quad
    \sup_{n\geq1}\operatorname{rank}(f_n-g_n)<\infty.
  \]
  By subadditivity of rank, this is an equivalence relation.
  For composable families of matrices, one has
  \[
    \operatorname{rank}(g_nf_n-g_n'f_n')
    =\operatorname{rank}\bigl( g_n(f_n-f_n')+(g_n-g_n')f_n'\bigr)\leq\operatorname{rank}(f_n-f_n')+\operatorname{rank}(g_n-g_n').
  \]
  Thus $\sim$ is compatible with composition and hence is a congruence on the category.
  Let $F:\mathcal C^{\mathbb N_{>0}} \longrightarrow \mathcal{C}^{\mathbb N_{>0}}/_\sim = \mathcal{D}$ be the corresponding quotient functor.
  For every $d\geq0$, take the projection $p_d=(I_d\ \ 0)$.
  Then $p_dj_d=I_d$, $j_dp_d= \begin{psmallmatrix}I_d&0\\0&0\end{psmallmatrix}$, and $\operatorname{rank}(j_dp_d-I_{d+1})=1$.
  Take any morphism $s=(s_n)_{n\geq1}:Y\longrightarrow Z$ in $S^{\mathbb N_{>0}}$.
  Coordinatewise, set
  \[
    r_n=
    \begin{cases}
      s_n^{-1}, & s_n\text{ is an isomorphism}, \\
      p_d,      & s_n=j_d.
    \end{cases}
  \]
  Then $r=(r_n)_{n\geq1}:Z\longrightarrow Y$ satisfies $rs=\operatorname{Id}_Y$, and
  \[
    \sup_{n\geq1}\operatorname{rank}
    (s_nr_n-\operatorname{Id}_{Z_n})\leq1.
  \]
  It follows that $F(r)F(s)=\operatorname{Id}_{FY}$ and $F(s)F(r)=\operatorname{Id}_{FZ}$.
  Therefore $F$ sends every morphism in $S^{\mathbb N_{>0}}$ to an isomorphism.
  By the universal property of localization, $F$ induces a functor $\widetilde F: (\mathcal C^{\mathbb N_{>0}}) [(S^{\mathbb N_{>0}})^{-1}] \longrightarrow \mathcal{D}$.
  But
  \[
    \operatorname{rank}(u_n-v_n)
    =\operatorname{rank}(I_n)=n,
    \qquad
    \sup_{n\geq1}\operatorname{rank}(u_n-v_n)=\infty,
  \]
  so $F(u)\ne F(v)$, and hence $[u]\ne[v]$.
  Together with $\Theta[u]=\Theta[v]$, $\Theta$ is not faithful.
\end{proof}


\begin{theorem}\label{cor:fraction-calculus-preservation}
  Let $\mathcal{A}$ be a category, assume that the class of morphisms $S \subseteq \mathsf{Mor}(\mathcal{A})$ contains all identity morphisms, and let $Q : \mathcal{A} \longrightarrow \mathcal{A}[S^{-1}]$ be the localization.
  Suppose that the category $\mathcal{A}$ and the class of morphisms $S$ satisfy the following two conditions:
  \begin{align*}
    N = \sup \big\{\ell(\mathfrak P) \mid \mathfrak P\in \langle\mathsf{Mor}(\mathcal A)\sqcup\overline{S}\rangle\big\} < \infty, \\
    \sup\big\{ \delta(\mathfrak P,\mathfrak Q) \mid [\mathfrak P]=[\mathfrak Q], \ \mathfrak P, \mathfrak Q\in \langle\mathsf{Mor}(\mathcal A)\sqcup\overline{S}\rangle, \ |\mathfrak P|, |\mathfrak Q| \leq N+1\big\} < \infty.
  \end{align*}
  Then the following statements hold:
  \begin{enumerate}[label=\textup{(\arabic*)}]
    \item If $\mathcal{A}$ has $I$-indexed coproducts and $S$ is closed under $I$-indexed coproducts, then $\mathcal{A}[S^{-1}]$ also has $I$-indexed coproducts and $Q$ preserves them;
    \item If $\mathcal{A}$ has $I$-indexed products and $S$ is closed under $I$-indexed products, then $\mathcal{A}[S^{-1}]$ also has $I$-indexed products and $Q$ preserves them.
  \end{enumerate}
\end{theorem}

\begin{proof}
  The proof is the same as that of Proposition~\ref{prop:finite-structure-localization}.
  It suffices to consider the coproduct case.
  If $\mathcal A$ is empty, the assertion is immediate, so assume that $\mathcal A$ is nonempty.
  Assume that $\mathcal{A}$ satisfies the two conditions above, that $\mathcal{A}$ has $I$-indexed coproducts, and that $S$ is closed under $I$-indexed coproducts.
  The coproduct functor $\coprod_I : \mathcal{A}^I \longrightarrow \mathcal{A}$ and the diagonal functor $\Delta^I : \mathcal{A} \longrightarrow \mathcal{A} ^I$ are respectively left and right adjoint.
  In particular, the right adjoint $\Delta^I$ sends $S$ to $S^I$; because $S$ is closed under $I$-indexed coproducts, the left adjoint $\coprod_I$ sends $S^I$ to $S$.
  By \cite[Lemma~1.1.6, p.~9]{KrauseHomologicalTheory}, or by the primary source \cite[Chapter~I, \S~1, Lemma~1.2, p.~7]{GabrielZisman}, the adjunction $\coprod_I \dashv \Delta^I$ therefore induces an adjunction on the localized categories
  \[
    \xymatrix@C=30mm@R=10mm{
    \mathcal A^I
    \ar@<0.8ex>[r]^{\coprod_I}
    \ar[d]_{Q^I}
    &
    \mathcal A
    \ar@<0.8ex>[l]^{\Delta^I}
    \ar[d]^{Q}
    \\
    (\mathcal A[S^{-1}])^I
    \ar@<0.8ex>[r]^-{\widehat {\coprod_I}}
    &
    \mathcal A[S^{-1}]
    \ar@<0.8ex>[l]^-{\widehat{\Delta^I}}
    } \qquad \xymatrix@R=10mm{
    Q \circ \coprod_I = \widehat{\coprod_I} \circ Q^I,\\
    Q^I \circ \Delta^I = \widehat{\Delta^I} \circ Q.
    }
  \]
  By the hypothesis, Theorem~\ref{thm:product-canonical-fullness}, and Proposition~\ref{prop:faithfulness-under-fullness}, the canonical functor $(\mathcal{A}^I)[(S^I)^{-1}] \longrightarrow (\mathcal{A}[S^{-1}])^I$ is an isomorphism of categories, so one may regard $Q^I$ as the localization functor of $\mathcal A^I$ at $S^I$.
  It is clear that $\widehat{\Delta^I}$ is still the diagonal functor, and $\widehat{\coprod_I}$ is its left adjoint.
  Therefore $\mathcal A[S^{-1}]$ has $I$-indexed coproducts, and $Q\circ\coprod_I=\widehat{\coprod_I}\circ Q^I$ shows that $Q$ preserves them.
\end{proof}


\section{Two applications for discrete indexing categories}\label{sec:discrete-applications}

This section gives two classes of sufficient conditions under which the canonical functor $\Theta_I$ is an equivalence of categories.
Throughout this section, $I$ is a set.

\subsection{calculus of fractions}

Gabriel and Zisman give the axioms for a calculus of fractions, fraction representations of localized morphisms, and a criterion for equality of two fractions in Chapter~I, \S\S~2.2--2.4 \cite[Chapter~I, \S\S~2.2--2.4, pp.~12--14]{GabrielZisman}.

\begin{definition}
  Let $\mathcal{A}$ be a category and $S \subseteq \mathsf{Mor}(\mathcal{A})$ a class of morphisms.
  The class $S$ is said to admit a right calculus of fractions if it contains all identity morphisms, is closed under composition, and satisfies the following conditions:
  \begin{enumerate}[label=\textup{(\arabic*)}]
    \item for $f:X\longrightarrow Y$ and $s:Y'\longrightarrow Y$, if $s\in S$, then there is a commutative square of the following form, with $t\in S$:
          \[
            \xymatrix@C=12mm@R=6mm{
            X'\ar[r]^{f'}\ar[d]_t&Y'\ar[d]^s\\
            X\ar[r]_f&Y.
            }
          \]
    \item for morphisms $f,g:X\longrightarrow Y$, if there is an $s\in S$ such that $s\circ f=s\circ g$, then there is a $t\in S$ such that $f\circ t=g\circ t$.
  \end{enumerate}
  A left calculus of fractions is obtained by dualizing these conditions.
\end{definition}

\begin{corollary}\label{thm:fraction-calculus-application}
  Let $\mathcal A$ be a category, $S\subseteq\mathsf{Mor}(\mathcal A)$ a class of morphisms, and $I$ a set.
  Let
  \[
    Q:\mathcal A\longrightarrow\mathcal A[S^{-1}],
    \qquad
    Q_I:\mathcal A^I\longrightarrow
    (\mathcal A^I)[(S^I)^{-1}].
  \]
  be the localizations.
  If $S$ admits a right calculus of fractions in $\mathcal{A}$, respectively a left calculus of fractions, then $S^I$ admits a right calculus of fractions in $\mathcal A^I$, respectively a left calculus of fractions, and the canonical functor obtained by factoring $Q^I$ through $Q_I$ is an isomorphism:
  \[
    \Theta_I:
    (\mathcal A^I)[(S^I)^{-1}]
    \longrightarrow(\mathcal A[S^{-1}])^I,
    \qquad
    \Theta_I\circ Q_I=Q^I.
  \]

  Assume that $S$ admits a left or right calculus of fractions.
  Furthermore, if $\mathcal A$ has $I$-indexed (co)products and $S$ is closed under $I$-indexed (co)products, then $\mathcal A[S^{-1}]$ has $I$-indexed (co)products and $Q$ preserves them.
\end{corollary}

\begin{proof}
  It suffices to prove the case of a right calculus of fractions; the left calculus case is dual.
  The axioms for a calculus of fractions can be verified coordinatewise, so $S^I$ clearly admits a right calculus of fractions.

  \medskip
  By the right calculus of fractions, every $\mathfrak P\in\langle\mathsf{Mor}(\mathcal A)\sqcup\overline S\rangle$ is equivalent to $f\ast\overline s$, where $s\in S$ \cite[Chapter~I, \S~2.4, p.~14]{GabrielZisman}.
  Taking $F=\varnothing$ and $N=2$ in Theorem~\ref{thm:product-canonical-fullness}, $\Theta_I$ is full.

  \medskip

  By Proposition~\ref{prop:faithfulness-under-fullness}, it remains only to verify that $d_i(3)$ has a uniform upper bound in order to show that $\Theta_I$ is faithful.
  First consider the rewritings needed to refine a fraction.
  If $s:Z\longrightarrow X$, $f:Z\longrightarrow Y$, $a:T\longrightarrow Z$, and $s,s\circ a\in S$, then
  \[
    f\ast\overline s
    \xlongleftrightarrow{\leq2}
    f\ast\overline s\ast(s\circ a)\ast\overline{s\circ a}\xlongleftrightarrow{\leq4}
    (f\circ a)\ast\overline{s\circ a}.
  \]
  Thus every such refinement uses at most $6$ elementary rewriting steps.
  If $[f\ast\overline s]=[g\ast\overline t]$, then the equality criterion for right fractions \cite[Chapter~I, \S\S~2.3--2.4, pp.~13--14]{GabrielZisman} gives $a,b$ such that
  \[
    s\circ a=t\circ b\in S,
    \qquad f\circ a=g\circ b.
  \]
  Each of the two right fractions reaches this common refinement in at most $6$ steps, so the number of elementary rewriting steps needed to pass to a common denominator is bounded by $\delta(f\ast\overline s,g\ast\overline t)\leq12$.

  A single arrow can be rewritten as a right fraction in at most $3$ elementary steps.
  Adding a forward arrow to the left of $f\ast\overline s$ requires only one composition; when adding $\overline t$, the axioms for right fractions give $a\in S$ and $h$ such that $f\circ a=t\circ h$, and hence
  \[
    \overline t\ast f\ast\overline s
    \xlongleftrightarrow{\leq6}
    \overline t\ast(f\circ a)\ast\overline{s\circ a}=\overline t\ast(t\circ h)\ast\overline{s\circ a}
    \xlongleftrightarrow{\leq3}
    h\ast\overline{s\circ a}.
  \]
  Therefore, a composition of at most $3$ arrows can be rewritten as a right fraction in at most $3+2\cdot9=21$ steps.
  For arbitrary $\mathfrak P,\mathfrak Q\in \langle\mathsf{Mor}(\mathcal A)\sqcup\overline S\rangle$, one has
  \[
    \begin{gathered}
      [\mathfrak P]=[\mathfrak Q],\quad
      |\mathfrak P|,|\mathfrak Q|\leq3
      \quad\Longrightarrow\quad
      \delta(\mathfrak P,\mathfrak Q)\leq21+12+21=54.
    \end{gathered}
  \]
  Thus $d_i(3)\leq54$ for all $i\in I$.
  Taking $N_0=2$, $F_0=F=\varnothing$, and $D=54$ in Proposition~\ref{prop:faithfulness-under-fullness}, one concludes that $\Theta_I$ is faithful.
  Since $\Theta_I$ is bijective on objects, it is an isomorphism of categories.

  \medskip

  Finally, $\ell(\mathfrak P)\leq2$ and $d_i(3)\leq54$ verify the two conditions in Theorem~\ref{cor:fraction-calculus-preservation}.
  Hence, whenever the corresponding $I$-indexed (co)products exist in $\mathcal A$ and $S$ is closed under them, they exist in $\mathcal A[S^{-1}]$ and are preserved by $Q$.
\end{proof}

\subsection{Model structures}

The term model structure below is used in the sense of Definition~\ref{def:model-structure}.
A model structure on a category $\mathcal A$ is determined by three classes of morphisms $(\mathsf{Cofib},\mathsf{Fib},\mathsf{Weq})$ containing all identity morphisms and closed under composition; their elements are called cofibrations, fibrations, and weak equivalences, respectively.
Morphisms in $\mathsf{TCofib} = \mathsf{Cofib}\cap\mathsf{Weq}$ are called trivial cofibrations, and morphisms in $\mathsf{TFib} = \mathsf{Fib}\cap\mathsf{Weq}$ are called trivial fibrations.
The following axioms for $(\mathsf{Cofib}, \mathsf{Fib}, \mathsf{Weq})$ come from Quillen's definition of a closed model category \cite[Chapter~I, \S~1]{Quillen}; see also \cite[Section~3, Definition~3.3]{DwyerSpalinski}.

\begin{definition}\label{def:model-structure}
  A triple $(\mathsf{Cofib},\mathsf{Fib},\mathsf{Weq})$ is called a model structure on a category $\mathcal A$ if the following conditions hold.
  \begin{enumerate}
    \item[CM1] For every composable pair of morphisms $g \circ f$, if two of $(f,g,g \circ f)$ are weak equivalences, then so is the third.
    \item[CM2] The three classes of morphisms $(\mathsf{Cofib},\mathsf{Fib},\mathsf{Weq})$ are closed under retracts.
          A morphism $f' : X' \longrightarrow Y'$ is called a retract of $f : X \longrightarrow Y$ in the arrow category if there exist morphisms represented by the dotted arrows that make the following diagram commute:
          \[
            \xymatrix@R=15pt{
            X' \ar@{.>}[r] \ar@{->}[d]_-{f'} \ar@/^/@{=}[rr]^{\operatorname{Id}_{X'}} & X \ar@{.>}[r] \ar@{->}[d]_-{f} & X' \ar@{->}[d]^-{f'} \\
            Y' \ar@{.>}[r] \ar@/_/@{=}[rr]_{\operatorname{Id}_{Y'}} & Y \ar@{.>}[r] & Y'
            }
          \]
    \item[CM3] For every commutative square
          \[
            \xymatrix@C=14mm@R=7mm{
            X\ar[r]^a\ar[d]_i&E\ar[d]^p\\
            Y\ar[r]_b\ar@{..>}[ur]^h&B
            },
          \]
          if either $i\in\mathsf{Cofib}$ and $p \in \mathsf{TFib}$, or $i \in \mathsf{TCofib}$ and $p \in \mathsf{Fib}$, there exists $h$ such that $h\circ i=a$ and $p\circ h=b$.
    \item[CM4] Every morphism $f : X \longrightarrow Y$ has factorizations of the following two forms:
          \begin{enumerate}
            \item[(i)] there exist $i : X \longrightarrow Z \in \mathsf{Cofib}$ and $p : Z \longrightarrow Y \in \mathsf{TFib}$ such that $f = p \circ i$;
            \item[(ii)] there exist $i' : X \longrightarrow Z' \in \mathsf{TCofib}$ and $p' : Z' \longrightarrow Y \in \mathsf{Fib}$ such that $f = p' \circ i'$.
          \end{enumerate}
  \end{enumerate}
\end{definition}

Quillen's original definition also requires the category to have all finite limits and finite colimits \cite[Chapter~I, \S~1]{Quillen}.
Egger proved that the existence of finite products and finite coproducts is sufficient for constructing the homotopy category \cite[p.~1 and Theorem~3.2]{EggerModelCategories}.
In what follows, $\mathcal A$ is assumed to have finite products and finite coproducts.
In particular, $\mathcal A$ has an initial object $\bot$ and a terminal object $\top$.

\begin{definition}
  Let the full subcategories of cofibrant objects, fibrant objects, trivially cofibrant objects, and trivially fibrant objects be denoted, respectively, by
  \[
    \begin{aligned}
      \mathcal C
       & =\{X\in\mathcal A\mid \bot\longrightarrow X\in\mathsf{Cofib}\},  &
      \mathcal F
       & =\{X\in\mathcal A\mid X\longrightarrow\top\in\mathsf{Fib}\},       \\
      \mathsf T\mathcal C
       & =\{X\in\mathcal A\mid \bot\longrightarrow X\in\mathsf{TCofib}\}, &
      \mathsf T\mathcal F
       & =\{X\in\mathcal A\mid X\longrightarrow\top\in\mathsf{TFib}\}.
    \end{aligned}
  \]
  For the definitions of cofibrant and fibrant objects, see \cite[Section~3, Remark~3.4]{DwyerSpalinski}; trivially cofibrant and trivially fibrant objects are defined in the same way.
\end{definition}

\begin{lemma}\label{lem:model-replacements}
  Let $\mathcal A$ carry a model structure.
  For every object $X\in\mathcal A$, there exist a cofibrant object $QX$ and a trivial fibration $p_X : QX \longrightarrow X$, as well as a fibrant object $RX$ and a trivial cofibration $i_X : X \longrightarrow RX$.
\end{lemma}

\begin{proof}
  Apply the factorization axioms in (CM4) to the morphisms $\bot\longrightarrow X$ and $X \longrightarrow \top$.
\end{proof}

For every object $X$, fix such factorizations once and for all, and consistently write
\[
  p_X:QX\longrightarrow X,\qquad i_X:X\longrightarrow RX
\]
for a cofibrant replacement and a fibrant replacement, respectively.
This notation follows \cite[Section~5, Lemmas~5.1 and~5.3, pp.~24--25]{DwyerSpalinski}.
Since the factorization axiom does not require the factorizations to be functorial, $Q$ and $R$ here are initially only choices on objects; for a given morphism, the required lift is chosen separately.

\medskip

Quillen defined the homotopy category as the Gabriel--Zisman localization at the weak equivalences and proved that it is equivalent to the category of homotopy classes between cofibrant-fibrant objects \cite[Chapter~I, \S~1, Theorem~1, pp.~24--27]{Quillen}.
Under the hypotheses of finite products and finite coproducts used in this paper, Egger later proved that the same conclusion remains valid \cite[Theorem~3.2]{EggerModelCategories}.

\begin{theorem}\label{thm:quillen-homotopy-category}
  Suppose that $\mathcal A$ has finite products and finite coproducts and carries a model structure $(\mathsf{Cofib},\mathsf{Fib},\mathsf{Weq})$.
  Let $\gamma:\mathcal A\longrightarrow \operatorname{Ho}(\mathcal A) =\mathcal A[\mathsf{Weq}^{-1}]$ be the localization.
  Then, for every $X \in \mathcal{C}$ and $Y \in \mathcal{F}$, every morphism in $\mathrm{Hom}_{\mathsf{Ho}(\mathcal{A})}(X,Y)$ has the form $\gamma (f)$ for some $f \in \mathrm{Hom}_\mathcal{A}(X,Y)$.
  Moreover, two morphisms $f,g:X\longrightarrow Y$ are equal in $\operatorname{Ho}(\mathcal A)$ if and only if they are homotopic.
\end{theorem}

\begin{proof}
  This is Quillen's fundamental theorem on model structures \cite[Chapter~I, \S~1, Theorem~1 and Corollary~1, pp.~24--27]{Quillen}; here the weakened version that requires only the existence of finite products and finite coproducts is used \cite[Theorem~3.2]{EggerModelCategories}.
  Egger's comparison is induced by the inclusion of the cofibrant-fibrant objects followed by the localization functor.
  Its full faithfulness therefore identifies the Hom set in the localization with the set of homotopy classes of morphisms from $X$ to $Y$, which gives both assertions.
\end{proof}

\begin{lemma}\label{lem:product-model-structure}
  Let $I$ be an indexing set, and let $\bigl(\mathcal A_i,(\mathsf{Cofib}_i,\mathsf{Fib}_i,\mathsf{Weq}_i)\bigr)_{i\in I}$ be a family of model categories with finite products and finite coproducts.
  Then the product category $\prod_{i\in I}\mathcal A_i$ has the coordinatewise model structure
  \[
    \left(\prod_{i\in I}\mathsf{Cofib}_i,\,
    \prod_{i\in I}\mathsf{Fib}_i,\,
    \prod_{i\in I}\mathsf{Weq}_i\right).
  \]
  This model structure is called the product model structure \cite[Example~1.1.6, p.~4]{HoveyModelCategories}.
\end{lemma}

For ordinary model categories, the following comparison equivalence for discrete indexing categories is recorded in \cite[Section~2.5, equation~(2.7) and footnote~14, pp.~41--42]{RiehlHomotopicalCategories}.
The proof below also applies under the weakened finite-(co)product hypotheses of Definition~\ref{def:model-structure}.

\begin{theorem}\label{thm:model-category-application}
  Let $\mathcal A$ be a model category and $I$ a discrete indexing set.
  Let
  \[
    \gamma:\mathcal A\longrightarrow\operatorname{Ho}(\mathcal A),
    \qquad
    \gamma_I:\mathcal A^I\longrightarrow
    \operatorname{Ho}(\mathcal A^I)
    =(\mathcal A^I)[(\mathsf{Weq}^I)^{-1}].
  \]
  be the localizations.
  The canonical functor obtained by factoring $\gamma^I$ through $\gamma_I$,
  \[
    \Theta_I:
    \operatorname{Ho}(\mathcal A^I)
    \longrightarrow(\operatorname{Ho}(\mathcal A))^I,
    \qquad
    \Theta_I\circ\gamma_I=\gamma^I
  \]
  is an isomorphism of categories.
  Equivalently, for any two families of objects $X=(X_i)_{i\in I}$ and $Y=(Y_i)_{i\in I}$, there is a canonical bijection
  \[
    \operatorname{Hom}_{\operatorname{Ho}(\mathcal A^I)}
    (\gamma_IX,\gamma_IY)
    \xlongrightarrow{\sim}
    \prod_{i\in I}
    \operatorname{Hom}_{\operatorname{Ho}(\mathcal A)}
    (\gamma X_i,\gamma Y_i).
  \]
\end{theorem}

\begin{proof}
  Take a composition of arrows $\mathfrak P$ from $X$ to $Y$ in $\langle\mathsf{Mor}(\mathcal A)\sqcup\overline{\mathsf{Weq}}\rangle$.
  By the cofibrant replacement, the fibrant replacement, and Theorem~\ref{thm:quillen-homotopy-category}, there is an $h:QX\longrightarrow RY$ such that $\gamma(h)=\gamma(i_Y)\circ[\mathfrak P]\circ\gamma(p_X)$.
  Therefore $\mathfrak P$ is equivalent to the following composition of three arrows:
  \[
    \xymatrix@C=14mm{
    X&
    QX\ar[l]_{p_X}\ar[r]^h&
    RY&
    Y\ar[l]_{i_Y}
    },
    \qquad
  [\mathfrak P]=[\overline{i_Y}\ast h\ast\overline{p_X}],
    \qquad
    \ell(\mathfrak P)\leq3.
  \]
  Taking $F=\varnothing$ and $N=3$ in Theorem~\ref{thm:product-canonical-fullness}, one concludes that $\Theta_I$ is full.

  \medskip

  By Proposition~\ref{prop:faithfulness-under-fullness}, it remains only to verify that $d_i(4)$ has a uniform upper bound.
  The notation $\xlongleftrightarrow{\leq m}$ continues to denote a sequence of at most $m$ elementary rewriting steps from \eqref{eq:gz-arrow-composition-relations}.
  First let $C \in \mathcal{C}$ be cofibrant, let $T \in \mathcal{F}$ be fibrant, and let $h_0,h_1:C\longrightarrow T$ satisfy $\gamma(h_0)=\gamma(h_1)$.
  By Theorem~\ref{thm:quillen-homotopy-category} and \cite[Lemma~4.6]{DwyerSpalinski}, there is a commutative diagram
  \[
    \xymatrix@C=12mm@R=7mm{
    C\amalg C\ar[r]^{(h_0, h_1)}\ar[d]_{(\operatorname{Id}_C, \operatorname{Id}_C)} \ar@{..>}[dr]^{(i_0, i_1)} &T\\
    C&\operatorname{Cyl}(C)\ar[u]_H\ar[l]_p
    }
  \]
  where $(i_0,i_1)\in\mathsf{Cofib}$ and $p\in \mathsf{TFib}$.
  Hence
  \[
    i_\varepsilon
    \xlongleftrightarrow{\leq2}
    \overline p\ast p\ast i_\varepsilon
    \xlongleftrightarrow{\leq2}\overline p,
    \qquad
    h_\varepsilon
    \longleftrightarrow H\ast i_\varepsilon
    \xlongleftrightarrow{\leq4}H\ast\overline p.
  \]
  Therefore
  \begin{equation}\label{eq:model-homotopy-rewriting}
    \delta(h_0,h_1)\leq10.
  \end{equation}

  \medskip

  Next take $C,D \in \mathcal{C}$, a weak equivalence $s:C\longrightarrow D$, an object $T \in \mathcal{F}$, and a morphism $h:C\longrightarrow T$.
  By (CM4) and (CM1), choose
  \[
    s=p\circ u,\qquad C\xlongrightarrow{u}Z\xlongrightarrow{p}D,\qquad u\in \mathsf{TCofib},\quad p\in \mathsf{TFib}.
  \]
  The following two lifting squares give $v:Z\longrightarrow T$ and $r:D\longrightarrow Z$:
  \[
    \xymatrix@C=12mm@R=7mm{
    C\ar[r]^h\ar[d]_u&T\ar[d]\\
    Z\ar[r]\ar@{..>}[ur]^v&\top
    }
    \qquad
    \xymatrix@C=12mm@R=7mm{
    \bot\ar[r]\ar[d]&Z\ar[d]^p\\
    D\ar[r]_{\operatorname{Id}_D}\ar@{..>}[ur]^r&D
    },
    \qquad
    v\circ u=h,\quad p\circ r=\operatorname{Id}_D.
  \]
  Therefore the following rewritings hold:
  \begin{equation}\label{eq:model-weak-inverse-rewriting}
    h\ast\overline s
    =(v\circ u)\ast\overline{p\circ u}
    \xlongleftrightarrow{\leq6}
    v\ast\overline p
    \xlongleftrightarrow{\leq5}v\circ r.
  \end{equation}

  \medskip

  For every morphism $f:U\longrightarrow V$, the lifting axiom gives a morphism $\widetilde f:QU\longrightarrow QV$:
  \[
    \xymatrix@C=16mm@R=7mm{
    \bot\ar[r]\ar[d]&QV\ar[d]^{p_V}\\
    QU\ar[r]_{f\circ p_U}
    \ar@{..>}[ur]^{\widetilde f}&V
    },
    \qquad
    p_V\circ\widetilde f=f\circ p_U.
  \]
  This commutative square gives the first rewriting below for every $f\in\mathsf{Mor}(\mathcal A)$.
  If in addition $f\in\mathsf{Weq}$, then $\widetilde f\in\mathsf{Weq}$, and let $\overline f:V\longrightarrow U$ and $\overline{\widetilde f}:QV\longrightarrow QU$ denote the inverse arrows in the localization corresponding to $f$ and $\widetilde f$, respectively; in this case one also has the second rewriting:
  \begin{equation}\label{eq:model-cofibrant-arrow-rewriting}
    f\xlongleftrightarrow{\leq4} p_V\ast\widetilde f\ast\overline{p_U}
    \quad\bigl(f\in\mathsf{Mor}(\mathcal A)\bigr), \qquad
    \overline f\xlongleftrightarrow{\leq8}
    p_U\ast\overline{\widetilde f}\ast\overline{p_V}
    \quad\bigl(f\in\mathsf{Weq}\bigr).
  \end{equation}

  Now take $\mathfrak P=a_m\ast\cdots\ast a_1 \in\langle\mathsf{Mor}(\mathcal A)\sqcup\overline{\mathsf{Weq}}\rangle$, where $1\leq m\leq4$, $a_k:X_{k-1}\longrightarrow X_k$, $X_0=X$, and $X_m=Y$.
  Let $p_k=p_{X_k}$.
  Applying \eqref{eq:model-cofibrant-arrow-rewriting} to each $a_k$ gives
  \[
    \begin{aligned}
      \mathfrak P
       & \xlongleftrightarrow{\leq8m}
      (p_m\ast b_m\ast\overline{p_{m-1}})
      \ast\cdots\ast(p_1\ast b_1\ast\overline{p_0})   \\
       & \xlongleftrightarrow{\leq2(m-1)}
      p_m\ast b_m\ast\cdots\ast b_1\ast\overline{p_0} \\
       & \xlongleftrightarrow{\leq3}
      \overline{i_Y}\ast(i_Y\circ p_m)
      \ast b_m\ast\cdots\ast b_1\ast\overline{p_0}.
    \end{aligned}
  \]
  where $b_k$ is an arrow from $QX_{k-1}$ to $QX_k$: it is either a morphism in $\mathcal A$ or the inverse arrow corresponding to a weak equivalence.
  Starting with $(i_Y\circ p_m)\ast b_m$, process the arrows $b_k$ successively from left to right: compose once for a forward arrow, and apply \eqref{eq:model-weak-inverse-rewriting} for an inverse arrow.
  In at most $11m$ steps, this gives a morphism $h:QX\longrightarrow RY$.
  Therefore
  \begin{equation}\label{eq:model-three-arrow-rewriting}
    \delta\bigl(\mathfrak P,
    \overline{i_Y}\ast h\ast\overline{p_X}\bigr)
    \leq8m+2(m-1)+3+11m\leq85.
  \end{equation}
  If $\mathfrak Q\in \langle\mathsf{Mor}(\mathcal A)\sqcup\overline{\mathsf{Weq}}\rangle$ satisfies $[\mathfrak P]=[\mathfrak Q]$ and $|\mathfrak Q|\leq4$, use the same cofibrant replacement $p_X$ and fibrant replacement $i_Y$ for $\mathfrak Q$ to obtain a morphism $h':QX\longrightarrow RY$.
  Then
  \[
    [h]=\gamma(i_Y)\circ[\mathfrak P]\circ\gamma(p_X)
    =\gamma(i_Y)\circ[\mathfrak Q]\circ\gamma(p_X)=[h'].
  \]
  Combining \eqref{eq:model-homotopy-rewriting} and \eqref{eq:model-three-arrow-rewriting}, one obtains
  \[
    \delta(\mathfrak P,\mathfrak Q)\leq85+10+85=180,
    \qquad d_i(4)\leq180\quad(i\in I).
  \]
  Taking $N_0=3$, $F_0=F=\varnothing$, and $D=180$ in Proposition~\ref{prop:faithfulness-under-fullness}, one concludes that $\Theta_I$ is faithful.
  Since $\Theta_I$ is bijective on objects, it is an isomorphism of categories.
\end{proof}

\begin{corollary}\label{cor:model-category-preservation}
  Let $(\mathcal A, (\mathsf{Cofib}, \mathsf{Fib}, \mathsf{Weq}))$ be a model category.
  Let $\gamma:\mathcal A\longrightarrow\operatorname{Ho}(\mathcal A)$ be the localization.
  Fix a set $I$.
  \begin{enumerate}[label=\textup{(\arabic*)}]
    \item If $\mathcal A$ has $I$-indexed coproducts and $\mathsf{Weq}$ is closed under these coproducts, then $\operatorname{Ho}(\mathcal A)$ has $I$-indexed coproducts and $\gamma$ preserves them;
    \item If $\mathcal A$ has $I$-indexed products and $\mathsf{Weq}$ is closed under these products, then $\operatorname{Ho}(\mathcal A)$ has $I$-indexed products and $\gamma$ preserves them.
  \end{enumerate}
  In particular, if $\mathcal A$ has all coproducts (products) and $\mathsf{Weq}$ is closed under them, then $\gamma$ preserves all coproducts (products).
\end{corollary}

\begin{proof}
  Take $S=\mathsf{Weq}$ and $Q=\gamma$ in Theorem~\ref{cor:fraction-calculus-preservation}.
  The preceding proof gives $\ell(\mathfrak P)\leq3$ and $d_i(4)\leq180$, so the two conditions in Theorem~\ref{cor:fraction-calculus-preservation} hold, and the stated conclusions about $I$-indexed (co)products and their preservation follow.
\end{proof}

\section{Conjectures and further questions}\label{sec:conjectures-and-problems}
\begin{conjecture}\label{conj:nonautomatic-additivity}
  There exist an additive category $\mathcal A$ and a class of morphisms $S$ such that both $\mathcal A$ and $\mathcal A[S^{-1}]$ are additive categories, while the localization functor $Q:\mathcal A\longrightarrow\mathcal A[S^{-1}]$ is not an additive functor.
\end{conjecture}

There is in fact an example in which $\mathcal A$ is preadditive and the ordinary category $\mathcal A[S^{-1}]$ can be equipped with a preadditive structure for which $Q$ does not preserve addition of the Hom sets.

\begin{example}\label{ex:preadditive-nonadditive-localization}
  Let $\mathcal A$ be the one-object category associated with the algebra $\mathbb F_5\times\mathbb F_2$: the category $\mathcal A$ has a single object $*$ and $\operatorname{End}_{\mathcal A}(*)=\mathbb F_5\times\mathbb F_2$.
  Composition and addition of morphisms are given by multiplication and addition in this algebra, respectively.
  Let $e=(1,0)$ and take $S=\{e\}$.
  Then $\mathcal A[S^{-1}]$ can be equipped with a preadditive structure for which the localization functor $Q:\mathcal A\longrightarrow \mathcal{A}[S^{-1}]$ does not preserve addition.
\end{example}

\begin{proof}
  Let $\mathcal B$ be the one-object category associated with the algebra $\mathbb F_5$.
  Define a functor
  \[
    F:\mathcal A\longrightarrow\mathcal B,
    \qquad F(*)=*,\qquad F(a,b)=a^3.
  \]
  Since $(ac)^3=a^3c^3$, $F$ is a functor, and $F(e)=1$.
  The universal property of localization therefore gives a unique functor
  \[
    \widetilde F:\mathcal A[S^{-1}]\longrightarrow\mathcal B,
    \qquad
    \widetilde F\circ Q=F.
  \]
  By construction, $Q((1,0))$ is an isomorphism.
  It is straightforward to verify that the following assignment satisfies the identity and composition laws and therefore defines a functor:
  \[
    G:\mathcal B\longrightarrow\mathcal A[S^{-1}],
    \qquad
    G(*)=Q(*),\qquad G(a)=Q(a^3,0).
  \]
  Then $\widetilde F\bigl(G(a)\bigr) =F(a^3,0)=a^9=a$.
  Hence $\widetilde F\circ G=\operatorname{Id}_{\mathcal B}$.
  On the other hand, $Q(a,b)=Q(e)\circ Q(a,b)=Q(a,0)$, and therefore
  \[
    G\bigl(\widetilde F(Q(a,b))\bigr) =G(a^3)=Q(a^9,0)=Q(a,0)=Q(a,b).
  \]
  Thus $(G\circ\widetilde F)\circ Q=Q$.
  By uniqueness in the universal property of localization, $G\circ\widetilde F =\operatorname{Id}_{\mathcal A[S^{-1}]}$.
  Hence $\widetilde F$ is an isomorphism of categories.
  Transport the addition of $\mathcal B$ along this isomorphism: for morphisms $u,v$ in the unique Hom set of $\mathcal A[S^{-1}]$, define
  \[
    u+_{S}v
    =G\bigl(\widetilde F(u)+\widetilde F(v)\bigr).
  \]
  With this addition, $\mathcal A[S^{-1}]$ is preadditive and $\widetilde F$ is an additive isomorphism.

  \medskip

  It remains to prove that $Q$ is not additive.
  If $Q$ preserved addition, then $F=\widetilde F\circ Q$ would also preserve addition.
  But
  \[
    F(e+e)=F(2,0)=2^3=3
    \neq2=1+1=F(e)+F(e) \qquad \text{in $\mathbb F_5$}
  \]
  Thus $Q$ does not preserve addition.
\end{proof}

For a set $I$, the product category $\mathcal{A}^I$ is precisely the functor category $\mathrm{Funct}(I, \mathcal{A})$.
What happens for a general indexing category $J$?

\begin{problem}\label{prob:functor-localization}
Let $J$ be a category.
Let $S^J$ be the class of all natural transformations in $\mathcal A^J$ that belong pointwise to $S$.
When is the canonical functor induced by localization
\[
  \Theta_{J}:(\mathcal A^{J})[(S^{J})^{-1}]\longrightarrow(\mathcal A[S^{-1}])^{J}
\]
an equivalence of categories, and when is it essentially surjective, full, or faithful, respectively?
\end{problem}

This property already fails in general for arrow categories.

\begin{example}\label{ex:morphism-category-localization}
  Let $\{a,s\}^*$ be the free monoid generated by the letters $a$ and $s$, and let $\mathcal A$ be the corresponding one-object category.
  Let $S=\{s^n\mid n\geq0\}$, and let $[1]$ be the category generated by the arrow $0\longrightarrow1$, so that $\mathcal A^{[1]}$ is the arrow category of $\mathcal A$.
  Then the canonical functor induced by localization
  \[
    \Theta_{[1]}:
    (\mathcal A^{[1]})[(S^{[1]})^{-1}]
    \longrightarrow
    (\mathcal A[S^{-1}])^{[1]}
  \]
  is not essentially surjective and hence is not an equivalence of categories.
\end{example}

\begin{proof}
  The localization $\mathcal A[S^{-1}]$ still has a single object, and its endomorphism monoid is $M=\langle a,s,\overline s\mid s\overline s=1=\overline s s\rangle \cong\mathbb N*\mathbb Z$.

  Observe that $(\mathcal A^{[1]})[(S^{[1]})^{-1}]$ and $\mathcal A^{[1]}$ have the same objects, namely all morphisms of $\mathcal{A}$.
  Take any object $m\in\{a,s\}^*$ of $(\mathcal A^{[1]})[(S^{[1]})^{-1}]$; the functor $\Theta_{[1]}$ sends it to the object $[m]$.
  A functor that distinguishes $[a\overline s a]$ from $[m]$ will now be constructed.
  Take
  \[
    V=\mathbb Q^2, \qquad T=\begin{psmallmatrix}1&1\\0&1\end{psmallmatrix}, \quad P=\begin{psmallmatrix}0&0\\1&1\end{psmallmatrix}.
  \]
  Then for every $n\in\mathbb Z$ one has $T^n=\begin{psmallmatrix}1&n\\0&1\end{psmallmatrix}$ and $PT^nP=(n+1)P$.
  Define a functor
  \[
    F:\mathcal A\longrightarrow\operatorname{Vect}_{\mathbb Q}, \qquad F(*)=V,
    \qquad F(a)=P,
    \quad F(s)=T.
  \]
  Every $F(s^n)=T^n$ is invertible, so the universal property of localization induces a functor $H:\mathcal A[S^{-1}]\longrightarrow \operatorname{Vect}_{\mathbb Q}$.
  In particular,
  \[
    H([a])=P,
    \qquad H([s])=T,
    \qquad H([\overline s])=T^{-1}.
  \]
  Taking $n=-1$ gives $H([a\overline s a])=PT^{-1}P=0$.

  Take any $m\in\{a,s\}^*$.
  If $m=s^r$ with $r\geq0$, then $H([m])=T^r\neq0$.
  If $m$ contains $t\geq1$ copies of the letter $a$, then there exist $n_0,\ldots,n_t\in\mathbb N$ such that $m=s^{n_0}a s^{n_1}a\cdots a s^{n_t}$.
  Repeatedly applying $PT^{n_j}P=(n_j+1)P$ gives
  \[
    H([m])
    =T^{n_0}PT^{n_1}P\cdots PT^{n_t}
    =\bigg(\prod_{j=1}^{t-1}(n_j+1)\bigg)
    T^{n_0}PT^{n_t}\neq0.
  \]
  When $t=1$, the empty product here is $1$.
  Thus $H([m])$ is nonzero for every $m\in\{a,s\}^*$.
  If $[a\overline s a]\cong[m]$, applying the functor
  \[
    H^{[1]}:(\mathcal A[S^{-1}])^{[1]}
    \longrightarrow(\operatorname{Vect}_{\mathbb Q})^{[1]},
  \]
  to this isomorphism would give an isomorphism between the zero morphism $0:V\longrightarrow V$ and the nonzero morphism $H([m]):V\longrightarrow V$.
  This is a contradiction.
  Therefore $[a\overline s a]$ is not isomorphic to any object in the image of $\Theta_{[1]}$, so $\Theta_{[1]}$ is not essentially surjective.
\end{proof}

\section*{Use of AI}
The manuscript was completed several months before the AI-assisted checks described here.
GPT-5.5 sol was used to locate references and check mathematical statements.
GPT-6 astra was also used in an attempt to resolve Conjecture~7.1, but this attempt yielded no progress.

\begin{thebibliography}{AGV72}
  \bibitem[AGV72]{SGA4}M. Artin, A. Grothendieck and J.-L. Verdier,
  \emph{Th\'eorie des topos et cohomologie \'etale des sch\'emas, Tome~1},
  Lecture Notes in Mathematics 269, Springer-Verlag, Berlin--Heidelberg, 1972,
  \url{https://doi.org/10.1007/BFb0081551}.
  \bibitem[Cis10]{CisinskiDerivable}D.-C. Cisinski,
  Cat\'egories d\'erivables,
  \emph{Bulletin de la Soci\'et\'e Math\'ematique de France} \textbf{138} (2010), no.~3, 317--393,
  \url{https://doi.org/10.24033/bsmf.2592}.
  \bibitem[DS95]{DwyerSpalinski}W. G. Dwyer and J. Spali\'nski,
  Homotopy theories and model categories,
  in I. M. James (ed.), \emph{Handbook of Algebraic Topology},
  North-Holland, Amsterdam, 1995, 73--126,
  \url{https://doi.org/10.1016/B978-0-444-81779-2/50003-1}.
  \bibitem[Egg06]{EggerModelCategories}J. M. Egger,
  Quillen model categories without equalisers or coequalisers,
  arXiv:math/0609808 [math.CT], 2006,
  \url{https://arxiv.org/abs/math/0609808}.
  \bibitem[GZ67]{GabrielZisman}P. Gabriel and M. Zisman,
  \emph{Calculus of Fractions and Homotopy Theory},
  Ergebnisse der Mathematik und ihrer Grenzgebiete 35,
  Springer-Verlag, Berlin--Heidelberg--New York, 1967,
  \url{https://doi.org/10.1007/978-3-642-85844-4}.
  \bibitem[Hov99]{HoveyModelCategories}M. Hovey,
  \emph{Model Categories},
  Mathematical Surveys and Monographs 63,
  American Mathematical Society, Providence, RI, 1999,
  \url{https://doi.org/10.1090/surv/063}.
  \bibitem[Kra10]{KrauseLocalizationTheory}H. Krause,
  Localization theory for triangulated categories,
  in T. Holm, P. J{\o}rgensen and R. Rouquier (eds.), \emph{Triangulated Categories},
  London Mathematical Society Lecture Note Series 375,
  Cambridge University Press, Cambridge, 2010, 161--235,
  \url{https://doi.org/10.1017/CBO9781139107075.005}.
  \bibitem[Kra22]{KrauseHomologicalTheory}H. Krause,
  \emph{Homological Theory of Representations},
  Cambridge Studies in Advanced Mathematics 195,
  Cambridge University Press, Cambridge, 2022,
  \url{https://doi.org/10.1017/9781108979108}.
  \bibitem[Mal05]{MaltsiniotisGrothendieck}G. Maltsiniotis,
  \emph{La th\'eorie de l'homotopie de Grothendieck},
  Ast\'erisque 301, Soci\'et\'e Math\'ematique de France, Paris, 2005,
  \url{https://doi.org/10.24033/ast.689}.
  \bibitem[Nan80]{NandaUniverseFractions}S. Nanda,
  A note on the universe of a category of fractions,
  \emph{Canadian Mathematical Bulletin} \textbf{23} (1980), no.~4, 425--427,
  \url{https://doi.org/10.4153/CMB-1980-062-9}.
  \bibitem[Qui67]{Quillen}D. G. Quillen,
  \emph{Homotopical Algebra},
  Lecture Notes in Mathematics 43, Springer-Verlag, Berlin--Heidelberg, 1967,
  \url{https://doi.org/10.1007/BFb0097438}.
  \bibitem[Rie22]{RiehlHomotopicalCategories}E. Riehl,
  Homotopical categories: from model categories to $(\infty,1)$-categories,
  in A. J. Blumberg, T. Gerhardt and M. A. Hill (eds.),
  \emph{Stable Categories and Structured Ring Spectra},
  Mathematical Sciences Research Institute Publications 69,
  Cambridge University Press, Cambridge, 2022, 5--74,
  \url{https://doi.org/10.1017/9781009128957.002}.
  \bibitem[Ser03]{SerpeKInjective}C. Serp\'e,
  Resolution of unbounded complexes in Grothendieck categories,
  \emph{Journal of Pure and Applied Algebra} \textbf{177} (2003), no.~1, 103--112,
  \url{https://doi.org/10.1016/S0022-4049(02)00075-0}.
  \bibitem[Tro97]{TroninProducts}S. N. Tronin,
  Products in categories of fractions and universal inversion of homomorphisms,
  \emph{Sbornik: Mathematics} \textbf{188} (1997), no.~10, 1521--1541,
  \url{https://doi.org/10.1070/SM1997v188n10ABEH000266}.
  \bibitem[Ver96]{VerdierDerived}J.-L. Verdier,
  \emph{Des cat\'egories d\'eriv\'ees des cat\'egories ab\'eliennes},
  Ast\'erisque 239, Soci\'et\'e Math\'ematique de France, Paris, 1996,
  \url{https://doi.org/10.24033/ast.364}.
\end{thebibliography}
\end{document}
